% Options for packages loaded elsewhere
\PassOptionsToPackage{unicode}{hyperref}
\PassOptionsToPackage{hyphens}{url}
\PassOptionsToPackage{dvipsnames,svgnames,x11names}{xcolor}
\documentclass[
  english,
  11pt,
  a4paper,
]{article}
\usepackage{xcolor}
\usepackage[margin=25mm]{geometry}
\usepackage{amsmath,amssymb}
\setcounter{secnumdepth}{-\maxdimen} % remove section numbering
\usepackage{iftex}
\ifPDFTeX
  \usepackage[T1]{fontenc}
  \usepackage[utf8]{inputenc}
  \usepackage{textcomp} % provide euro and other symbols
\else % if luatex or xetex
  \usepackage{unicode-math} % this also loads fontspec
  \defaultfontfeatures{Scale=MatchLowercase}
  \defaultfontfeatures[\rmfamily]{Ligatures=TeX,Scale=1}
\fi
\usepackage{lmodern}
\ifPDFTeX\else
  % xetex/luatex font selection
\fi
% Use upquote if available, for straight quotes in verbatim environments
\IfFileExists{upquote.sty}{\usepackage{upquote}}{}
\IfFileExists{microtype.sty}{% use microtype if available
  \usepackage[]{microtype}
  \UseMicrotypeSet[protrusion]{basicmath} % disable protrusion for tt fonts
}{}
\makeatletter
\@ifundefined{KOMAClassName}{% if non-KOMA class
  \IfFileExists{parskip.sty}{%
    \usepackage{parskip}
  }{% else
    \setlength{\parindent}{0pt}
    \setlength{\parskip}{6pt plus 2pt minus 1pt}}
}{% if KOMA class
  \KOMAoptions{parskip=half}}
\makeatother
\usepackage{longtable,booktabs,array}
\newcounter{none} % for unnumbered tables
\usepackage{calc} % for calculating minipage widths
% Correct order of tables after \paragraph or \subparagraph
\usepackage{etoolbox}
\makeatletter
\patchcmd\longtable{\par}{\if@noskipsec\mbox{}\fi\par}{}{}
\makeatother
% Allow footnotes in longtable head/foot
\IfFileExists{footnotehyper.sty}{\usepackage{footnotehyper}}{\usepackage{footnote}}
\makesavenoteenv{longtable}
\ifLuaTeX
\usepackage[bidi=basic,shorthands=off]{babel}
\else
\usepackage[bidi=default,shorthands=off]{babel}
\fi
\ifLuaTeX
  \usepackage{selnolig} % disable illegal ligatures
\fi
\setlength{\emergencystretch}{3em} % prevent overfull lines
\providecommand{\tightlist}{%
  \setlength{\itemsep}{0pt}\setlength{\parskip}{0pt}}
\setlength{\emergencystretch}{3em}
\setlength{\parindent}{0pt}
\setlength{\parskip}{0.45em}
\sloppy
\usepackage{bookmark}
\IfFileExists{xurl.sty}{\usepackage{xurl}}{} % add URL line breaks if available
\urlstyle{same}
\hypersetup{
  pdftitle={Nonexistence of Simple Graphs with Laplacian Spectrum
\{0,1,\ldots,n-1\} for n=16 and 20: Exact Computer-Assisted Proofs},
  pdfauthor={Jihun Kim (Independent Researcher)},
  pdflang={en},
  pdfsubject={Exact certificate-based computer-assisted proofs in
spectral graph theory},
  pdfkeywords={spectral graph theory, Laplacian
spectrum, computer-assisted proof, Farkas certificate, exact
computation},
  colorlinks=true,
  linkcolor={blue},
  filecolor={Maroon},
  citecolor={Blue},
  urlcolor={blue},
  pdfcreator={LaTeX via pandoc}}

\title{Nonexistence of Simple Graphs with Laplacian Spectrum
\(\{0,1,\ldots,n-1\}\) for \(n=16\) and \(20\): Exact Computer-Assisted
Proofs}
\usepackage{etoolbox}
\makeatletter
\providecommand{\subtitle}[1]{% add subtitle to \maketitle
  \apptocmd{\@title}{\par {\large #1 \par}}{}{}
}
\makeatother
\subtitle{Preprint, version 0.1.0}
\author{Jihun Kim (Independent Researcher)}
\date{24 August 2026}

\begin{document}
\maketitle

\textbf{MSC 2020:} 05C50, 05C85, 15A18.

\begin{center}\rule{0.5\linewidth}{0.5pt}\end{center}

\subsection{Abstract}\label{abstract}

We establish the \(n=16\) and \(n=20\) cases of the \(S_{n,n}\)
conjecture. There is no simple graph on \(16\) vertices with Laplacian
spectrum \(\{0,1,\ldots,15\}\), and no simple graph on \(20\) vertices
with Laplacian spectrum \(\{0,1,\ldots,19\}\). Both proofs combine exact
enumeration of forced degree data with linear relaxations derived from
spectral moments, graph and complement identities, and degree-class
incidence constraints. The order-16 proof eliminates \(2{,}077\)
candidate degree sequences with \(2{,}143\) exact rational Farkas leaves
organized into degree-reduction certificates and 42 integer branch
trees. The order-20 proof enumerates \(160{,}244\)
degree-sequence/triangle-count pairs, reduces them to \(80{,}122\)
complement representatives, and eliminates them with \(80{,}124\) exact
integer Farkas certificates, including four exhaustive branch leaves.
All certificate checks use integer or rational arithmetic, and both
coverage ledgers have zero uncovered cases. These finite-order results
do not settle the conjecture in general.

\clearpage
\setcounter{tocdepth}{2}
\tableofcontents
\clearpage

\subsection{Introduction and relation to previous
work}\label{introduction-and-relation-to-previous-work}

For \(1\le i\le n\), Fallat, Kirkland, Molitierno and Neumann introduced

\[
S_{i,n}=\{0,1,\ldots,n\}\setminus\{i\}
\]

and asked which such sets occur as the Laplacian spectrum of a simple
graph {[}1{]}. Their \(S_{n,n}\) conjecture states that

\[
S_{n,n}=\{0,1,\ldots,n-1\}
\]

is not Laplacian-realizable for any \(n\ge2\). Fallat et al.~established
the conjecture for \(n\le11\), prime \(n\), and \(n\equiv2,3\pmod4\)
{[}1{]}. Das, Lee and Cheon derived additional diameter restrictions on
any realizing graph and its complement {[}2{]}. Goldberger and Neumann
proved the conjecture for all sufficiently large orders {[}3{]}. Hameed,
Khan and Tyaglov surveyed the remaining range and, among other
structural results, excluded Cartesian-product realizations {[}4{]}.
Johnston, Plosker and Varona subsequently established the unrestricted
order-12 case and described \(n=16\) as the smallest open order in July
2026 {[}5{]}.

The present paper treats exactly two orders, \(16\) and \(20\). Within
the finite cases and order classes covered by the cited results, they
are the two smallest orders not covered: 17 and 19 are prime, while
\(18\equiv2\pmod4\). The least order not covered by that cited
collection after the present two theorems is 21. The proofs share a
mathematical architecture but not a single certificate corpus. The
order-16 computation uses rational Farkas certificates and complete
integer branch trees for 42 final class models. The order-20 computation
uses a complement quotient followed by three continuous relaxations and
a final two-root integer split.

A literature search completed on 24 August 2026 found no earlier public
proof of either case.

\subsection{1. Main results and proof
architecture}\label{main-results-and-proof-architecture}

Let \(G\) be a finite simple undirected graph, let \(A\) be its
adjacency matrix, let \(D=\operatorname{diag}(d_v)\), and let \(L=D-A\)
be its Laplacian.

\subsubsection{Theorem 1 (order 16)}\label{theorem-1-order-16}

There is no simple undirected graph \(G\) on \(16\) vertices such that

\[
\operatorname{spec}L(G)=\{0,1,2,\ldots,15\}.
\]

\subsubsection{Theorem 2 (order 20)}\label{theorem-2-order-20}

There is no simple undirected graph \(G\) on \(20\) vertices such that

\[
\operatorname{spec}L(G)=\{0,1,2,\ldots,19\}.
\]

Each proof has four logical components:

\begin{enumerate}
\def\labelenumi{\arabic{enumi}.}
\tightlist
\item
  derive necessary spectral, degree and triangle conditions;
\item
  enumerate every discrete candidate satisfying those conditions;
\item
  map every hypothetical target graph to a feasible point of a sequence
  of necessary linear or mixed-integer relaxations;
\item
  verify exact Farkas certificates, together with exhaustive integer
  branch splits where required, that eliminate every candidate.
\end{enumerate}

Floating-point optimization was used only during discovery. The proofs
rely on the analytic derivations, exact enumeration, certificate
verification and coverage ledgers described below.

Appendix A gives the complete enumeration and coverage ledgers.
Appendices B and C give the variables, rows, bounds, serialization and
certificate semantics used by the order-16 and order-20 verifiers.

\subsection{2. Shared spectral and graph
identities}\label{shared-spectral-and-graph-identities}

Assume temporarily that \(G\) has \(n\) vertices and Laplacian spectrum
\(\{0,1,\ldots,n-1\}\). The zero eigenvalue is simple, so \(G\) is
connected. The first three spectral power sums give

\[
\sum_v d_v=\operatorname{tr}L=\frac{n(n-1)}2,
\tag{2.1}
\]

\[
\sum_v d_v^2
=\operatorname{tr}L^2-\operatorname{tr}L
=\frac{n(n-1)(n-2)}3,
\tag{2.2}
\]

and

\[
\sum_v d_v^3-6t(G)
=\operatorname{tr}L^3-3\sum_vd_v^2.
\tag{2.3}
\]

Here we used

\[
(L^2)_{vv}=d_v^2+d_v,
\qquad
\operatorname{tr}L^3=\sum_vd_v^3+3\sum_vd_v^2-6t(G).
\]

The complement satisfies

\[
L(\overline G)=nI-J-L(G).
\]

On \(\mathbf1^\perp\), eigenvalue \(k\) is sent to \(n-k\). Hence the
target spectrum is complement-invariant. Direct counting also gives

\[
t(G)+t(\overline G)
=\binom n3-\frac12\sum_vd_v(n-1-d_v)
=\frac{n(n-1)(n-5)}{12}.
\tag{2.4}
\]

For a local degree bound, choose an orthonormal Laplacian eigenbasis
\(u_k\) and put \(w_{v,k}=u_k(v)^2\). Connectedness gives
\(w_{v,0}=1/n\), while

\[
\sum_kkw_{v,k}=d_v,
\qquad
\sum_kk^2w_{v,k}=d_v^2+d_v.
\tag{2.5}
\]

For \(1\le k\le n-1\),

\[
k^2\le nk-(n-1)
\]

because \((k-1)(n-1-k)\ge0\). Substitution in (2.5) gives

\[
nd_v^2-n(n-1)d_v+(n-1)^2\le0.
\tag{2.6}
\]

This proves the exact integer ranges used below:

\[
2\le d_v\le13\quad(n=16),
\qquad
2\le d_v\le17\quad(n=20).
\tag{2.7}
\]

Finally, every triangle contributes at each of its three vertices, so

\[
3t(G)\le\sum_v\binom{d_v}{2}.
\]

Combining this bound for \(G\) and \(\overline G\) with (2.4) yields

\[
54\le t(G)\le166\quad(n=16),
\qquad
127\le t(G)\le348\quad(n=20).
\tag{2.8}
\]

These bounds are consequences of the target spectrum, not imported
pruning assumptions.

\subsection{3. Common relaxation
variables}\label{common-relaxation-variables}

Let \(m_d\) be the multiplicity of degree \(d\). For an order-\(n\)
candidate define the nonnegative real spectral weights

\[
Y_{d,k}=n\sum_{v:d_v=d}u_k(v)^2,
\qquad 1\le k\le n-1.
\tag{3.1}
\]

Every target graph satisfies

\[
\sum_kY_{d,k}=(n-1)m_d,
\quad
\sum_kkY_{d,k}=nm_dd,
\quad
\sum_kk^2Y_{d,k}=nm_d(d^2+d),
\tag{3.2}
\]

and, for each \(k\),

\[
\sum_dY_{d,k}=n.
\tag{3.3}
\]

For a vertex \(v\) of degree \(d\), put

\[
s_v=\sum_{u\sim v}d_u,
\qquad
q_v=\#\{\text{triangles containing }v\},
\qquad
h_v=s_v-2q_v.
\]

Direct multiplication gives

\[
(L^3)_{vv}=d^3+2d^2+h_v.
\tag{3.4}
\]

If \(R_v=V\setminus(N(v)\cup\{v\})\), then

\[
h_v=d+e(N(v),R_v),
\qquad d\le h_v\le d(n-d).
\tag{3.5}
\]

For a degree class define

\[
S_d=\sum_{v:d_v=d}s_v,
\qquad
H_d=\sum_{v:d_v=d}h_v.
\]

Equations (3.1) and (3.4) imply

\[
\sum_kk^3Y_{d,k}=n\{m_d(d^3+2d^2)+H_d\}.
\tag{3.6}
\]

Summing (3.5) over the degree-\(d\) class gives the two necessary
Stage-2 inequalities

\[
m_dd\le H_d\le m_dd(n-d).
\tag{3.7}
\]

In both serialized Stage-2 models, \(H_d\) is not a separate variable.
It is the affine abbreviation

\[
H_d=\frac1n\sum_k k^3Y_{d,k}-m_d(d^3+2d^2)
\]

obtained from (3.6). The variables \(S_d\) are allowed to be real in the
relaxation, even though the values supplied by an actual graph are
integers. After eliminating \(H_d\), (3.7) and the bounds below are
linear rows in \(Y,S\).

The neighbor-degree order bounds, triangle capacities,
complement-triangle identity and \(\sum_dS_d=\sum_vd_v^2\) complete the
Stage-2 relaxation. If \(\ell_d,u_d\) are the sums of the \(d\) smallest
and largest entries after removing one occurrence of \(d\), then

\[
m_d\ell_d\le S_d\le m_du_d,
\tag{3.8}
\]

while

\[
H_d\le S_d\le H_d+2m_d\binom d2.
\tag{3.9}
\]

Writing \(M=n(n-1)/4\) for the forced edge count, the local complement
formula

\[
\overline q_v=\binom{n-1-d}{2}-M+s_v-q_v
\tag{3.10}
\]

gives

\[
m_d\left(2M-2\binom{n-1-d}{2}\right)
\le S_d+H_d\le2Mm_d.
\tag{3.11}
\]

All these rows are necessary conditions. The models deliberately omit
several nonlinear graph-realizability constraints; infeasibility of a
relaxation is therefore sufficient for nonexistence, whereas feasibility
is not evidence of a graph.

\subsection{4. The order-16 certified
exhaustion}\label{the-order-16-certified-exhaustion}

\subsubsection{4.1 Forced data and exact
enumeration}\label{forced-data-and-exact-enumeration}

For \(n=16\), (2.1)--(2.3) specialize to

\[
\sum_vd_v=120,
\qquad
\sum_vd_v^2=1120,
\qquad
\sum_vd_v^3-6t=11040.
\tag{4.1}
\]

The complement triangle total is \(220\), the degree range is
\([2,13]\), and the triangle range is \([54,166]\). Exact
nondecreasing-sequence enumeration has the following ledger.

{\def\LTcaptype{none} % do not increment counter
\begin{longtable}[]{@{}
  >{\raggedright\arraybackslash}p{(\linewidth - 6\tabcolsep) * \real{0.2000}}
  >{\raggedleft\arraybackslash}p{(\linewidth - 6\tabcolsep) * \real{0.2667}}
  >{\raggedleft\arraybackslash}p{(\linewidth - 6\tabcolsep) * \real{0.2667}}
  >{\raggedleft\arraybackslash}p{(\linewidth - 6\tabcolsep) * \real{0.2667}}@{}}
\toprule\noalign{}
\begin{minipage}[b]{\linewidth}\raggedright
Order-16 layer
\end{minipage} & \begin{minipage}[b]{\linewidth}\raggedleft
Source
\end{minipage} & \begin{minipage}[b]{\linewidth}\raggedleft
Eliminated
\end{minipage} & \begin{minipage}[b]{\linewidth}\raggedleft
Survivors
\end{minipage} \\
\midrule\noalign{}
\endhead
\bottomrule\noalign{}
\endlastfoot
moment equations & -- & -- & \(2{,}249\) \\
Erdős--Gallai graphicality & \(2{,}249\) & \(16\) & \(2{,}233\) \\
integral triangle count and bounds & \(2{,}233\) & \(156\) &
\(2{,}077\) \\
Stage 1 spectral weights & \(2{,}077\) & \(1{,}359\) & \(718\) \\
Stage 2 local third moments & \(718\) & \(676\) & \(42\) \\
integer class branch trees & \(42\) & \(42\) & \(0\) \\
\end{longtable}
}

A second implementation enumerates the moment solutions by degree
multiplicities and checks graphicality by Havel-Hakimi instead of
reusing the main combinations/Erdős--Gallai route. It reproduces
\(2{,}249\), \(2{,}233\) and \(2{,}077\) exactly.

The recurrence, pruning invariant, graphicality test, triangle filter,
deterministic record order and exact ledger-partition obligations for
this enumeration are specified in Appendix A.

\subsubsection{4.2 Spectral and local-third-moment
leaves}\label{spectral-and-local-third-moment-leaves}

The Stage-1 system consists of (3.2)--(3.3) at \(n=16\). Exact rational
Farkas certificates eliminate \(1{,}359\) candidates. Stage 2 appends
the real variables \(S_d\), uses the aggregate bounds (3.7)--(3.9) and
(3.11) with \(H_d\) eliminated by the affine expression obtained from
(3.6), and adds

\[
\sum_dS_d=1120.
\]

Exact rational Farkas certificates eliminate another \(676\) candidates,
leaving 42 sequences.

\subsubsection{4.3 Integer degree-class
model}\label{integer-degree-class-model}

For a surviving degree multiset let \(D\) be its distinct degrees and
put

\[
P_{a,b}=\begin{cases}
\binom{m_a}{2},&a=b,\\
m_am_b,&a<b.
\end{cases}
\]

The model retains the real variables \(Y_{d,k}\) and adds integer graph
counts:

\begin{itemize}
\tightlist
\item
  \(E_{a,b}\), the number of graph edges with endpoint degree multiset
  \(\{a,b\}\);
\item
  \(T_{a,b,c}\), the number of graph triangles with degree multiset
  \(\{a,b,c\}\);
\item
  \(C_{a,b,c}\), the analogous complement-triangle count indexed by
  original degrees.
\end{itemize}

Their bounds are the elementary pair and triple capacities. Stub
equations, graph and complement triangle totals, classwise third-moment
equations, edge-triangle incidence bounds and degree-class wedge bounds
are imposed. For example,

\[
2E_{d,d}+\sum_{b\ne d}E_{\min(d,b),\max(d,b)}=m_dd,
\tag{4.2}
\]

\[
\sum T_{a,b,c}=t,
\qquad
\sum C_{a,b,c}=220-t,
\tag{4.3}
\]

and, with \(\nu_d(a,b,c)\) the multiplicity of \(d\) in a triple,

\[
\sum_kk^3Y_{d,k}
=16\left[m_d(d^3+2d^2)+S_d^G-2R_d^G\right],
\tag{4.4}
\]

where

\[
S_d^G=2dE_{d,d}+\sum_{b\ne d}bE_{\min(d,b),\max(d,b)},
\qquad
R_d^G=\sum_{a\le b\le c}\nu_d(a,b,c)T_{a,b,c}.
\]

The complement uses \(P_{a,b}-E_{a,b}\), original degree \(d\) maps to
\(15-d\), and spectral weights satisfy

\[
\overline Y_{15-d,k}=Y_{d,16-k}.
\tag{4.5}
\]

Any graph with the prescribed spectrum determines a feasible mixed
point: \(Y\) is generally real, whereas \(E,T,C\) are integer. The
converse is not asserted. The complete triple set, capacities, graph and
complement rows, wedge and incidence inequalities, variable order and
branch-tree serialization are given in Appendix B; that appendix is the
formal definition of the model whose certificates are checked here.

\subsubsection{4.4 Exact branch-and-Farkas
proof}\label{exact-branch-and-farkas-proof}

At an order-16 branch-tree node, the model has rational rows

\[
Ax\le b,
\qquad
Cx=e.
\]

A leaf stores rational multipliers \(y,z\) satisfying

\[
y\ge0,
\qquad
y^{\mathsf T}A+z^{\mathsf T}C=0,
\qquad
y^{\mathsf T}b+z^{\mathsf T}e=-1.
\tag{4.6}
\]

Any feasible \(x\) would yield \(0\le-1\), a contradiction. An internal
node branches on an explicitly integer variable \(x_j\) using

\[
x_j\le a
\qquad\text{or}\qquad
x_j\ge a+1.
\tag{4.7}
\]

The cases are disjoint and exhaustive for every integer value in the
parent interval. The 42 trees contain 174 nodes, 66 branch nodes and 108
exact Farkas leaves, with maximum depth 10. Together with the
\(1{,}359+676\) earlier leaves, the order-16 proof verifies \(2{,}143\)
exact Farkas certificates. No target candidate survives the certified
computation. Lemma 3 below supplies the graph-to-model implication that
completes the proof of Theorem 1.

\subsubsection{Computational Proposition 1 (order-16 exact
coverage)}\label{computational-proposition-1-order-16-exact-coverage}

The order-16 ledger partitions all \(2{,}077\) candidates into
\(1{,}359\) Stage-1 leaves, \(676\) Stage-2 leaves and \(42\) class
roots. The branch forest replaces those roots by \(108\) certified
leaves, so \(2{,}143\) rational Farkas certificates cover every
candidate.

\subsection{5. The order-20 certified
exhaustion}\label{the-order-20-certified-exhaustion}

\subsubsection{5.1 Forced data, enumeration and complement
quotient}\label{forced-data-enumeration-and-complement-quotient}

For \(n=20\), the shared identities specialize to

\[
\sum_vd_v=190,
\qquad
\sum_vd_v^2=2280,
\qquad
\sum_vd_v^3-6t=29260.
\tag{5.1}
\]

The complement triangle total is \(475\), the degree range is
\([2,17]\), and the triangle range is \([127,348]\). Exact enumeration
gives

\[
209{,}932\longrightarrow200{,}108\longrightarrow160{,}244
\tag{5.2}
\]

moment sequences, graphical sequences, and graphical
degree-sequence/triangle-count pairs. The primary enumerator lists
degree multiplicities and uses Erdős--Gallai. A
position-by-position/Havel--Hakimi audit and a degree-alphabet
meet-in-the-middle/Kleitman-Wang audit reproduce the same counts and
complete uncompressed candidate bytes.

The exact enumeration recurrences, pruning invariants, record
serialization, complement-orbit checks and ledger partitions are
specified in Appendix A.

For a pair \((s,t)\), where \(s=(d_1,\ldots,d_{20})\) is nondecreasing,
define

\[
\kappa(s,t)
=\left(\operatorname{sort}(19-d_1,\ldots,19-d_{20}),475-t\right).
\tag{5.3}
\]

The map is a fixed-point-free involution on the 160,244 candidates: a
fixed point would require \(2t=475\). Retaining the lexicographically
smaller member of each orbit leaves exactly \(80{,}122\)
representatives. If a target graph realizes either member, it or its
complement realizes the retained one.

\subsubsection{5.2 Three necessary relaxation
layers}\label{three-necessary-relaxation-layers}

Stage 1 is (3.2)--(3.3) at \(n=20\). Stage 2 appends the real variables
\(S_d\), uses the aggregate bounds (3.7)--(3.9) and (3.11) with \(H_d\)
eliminated by the affine expression obtained from (3.6), and adds
\(\sum_dS_d=2280\). Stage 3 uses the real spectral weights together with
degree-pair and degree-triple variables \(E_{a,b},T_z,C_z\). For a
nondecreasing triple \(z\) with degree multiplicities \(\nu_d(z)\), put

\[
P_z=\prod_d\binom{m_d}{\nu_d(z)}.
\]

Only triples with \(P_z>0\) are serialized. Define the pair multiplicity
in a triple by

\[
\pi_{a,b}(z)=
\begin{cases}
\binom{\nu_a(z)}2,&a=b,\\
\nu_a(z)\nu_b(z),&a<b.
\end{cases}
\tag{5.4a}
\]

The elementary capacities are

\[
0\le E_{a,b}\le P_{a,b},
\qquad
0\le T_z,C_z\le P_z.
\tag{5.4}
\]

The model imposes stub equations, graph and complement triangle totals,
graph and complement classwise \(L^3\) equations, pair-incidence bounds
and wedge bounds. For a pair \((a,b)\), graph triangles incident with
that pair satisfy

\[
\sum_z\pi_{a,b}(z)T_z
\le(\min(a,b)-1)E_{a,b},
\tag{5.5}
\]

while complement triangles satisfy

\[
\sum_z\pi_{a,b}(z)C_z
\le(18-\max(a,b))(P_{a,b}-E_{a,b}).
\tag{5.6}
\]

In the 343 ordinary Stage-3 leaves, \(E,T,C\) are allowed to be real.
Thus each certificate proves infeasibility of a continuous relaxation
containing every actual graph point. Integrality metadata is consulted
only for the final branch variable described below.

Appendix C states the full graph and complement classwise equations,
stub and triangle-total rows, pair-incidence and wedge inequalities,
standard-form conversion, deterministic row order and sparse-certificate
serialization reconstructed by the verifier.

\subsubsection{5.3 Exact standard-form
certificates}\label{exact-standard-form-certificates}

For each order-20 leaf, finite lower bounds are shifted away,
inequalities and finite upper bounds receive nonnegative slack
variables, and the verifier obtains

\[
Bw=h,
\qquad
w\ge0,
\tag{5.7}
\]

with integer \(B,h\). A stored integer multiplier \(z\) is a
contradiction if

\[
B^{\mathsf T}z\ge0,
\qquad
h^{\mathsf T}z<0.
\tag{5.8}
\]

Indeed, a feasible \(w\) would imply

\[
h^{\mathsf T}z=w^{\mathsf T}B^{\mathsf T}z\ge0.
\]

The dedicated slack columns force the multipliers of inequality and
upper-bound rows to be nonnegative. The verifier checks integer
arithmetic, row signs, row order and every sparse multiplier exactly.

\subsubsection{5.4 Coverage and final integer
split}\label{coverage-and-final-integer-split}

The exact order-20 coverage ledger is:

{\def\LTcaptype{none} % do not increment counter
\begin{longtable}[]{@{}
  >{\raggedright\arraybackslash}p{(\linewidth - 6\tabcolsep) * \real{0.2000}}
  >{\raggedleft\arraybackslash}p{(\linewidth - 6\tabcolsep) * \real{0.2667}}
  >{\raggedleft\arraybackslash}p{(\linewidth - 6\tabcolsep) * \real{0.2667}}
  >{\raggedleft\arraybackslash}p{(\linewidth - 6\tabcolsep) * \real{0.2667}}@{}}
\toprule\noalign{}
\begin{minipage}[b]{\linewidth}\raggedright
Order-20 layer
\end{minipage} & \begin{minipage}[b]{\linewidth}\raggedleft
Source
\end{minipage} & \begin{minipage}[b]{\linewidth}\raggedleft
Exact Farkas leaves
\end{minipage} & \begin{minipage}[b]{\linewidth}\raggedleft
Survivors
\end{minipage} \\
\midrule\noalign{}
\endhead
\bottomrule\noalign{}
\endlastfoot
complement representatives & \(80{,}122\) & -- & \(80{,}122\) \\
Stage 1 spectral weights & \(80{,}122\) & \(63{,}235\) & \(16{,}887\) \\
Stage 2 local third moments & \(16{,}887\) & \(16{,}542\) & \(345\) \\
Stage 3 class model & \(345\) & \(343\) & \(2\) roots \\
integer split & \(2\) roots & \(4\) & \(0\) \\
\end{longtable}
}

The two parent roots are

\[
s^{(0)}=(2,2,2,4,6,6,7,7,9,9,9,10,12,12,14,14,15,16,17,17),
\qquad t^{(0)}=229,
\tag{5.9}
\]

and

\[
s^{(1)}=(2,2,2,5,5,6,6,8,9,9,9,10,12,12,14,14,15,16,17,17),
\qquad t^{(1)}=231.
\tag{5.10}
\]

Both have exactly three degree-2 vertices. Variable \(E_{2,2}\) is an
actual integer graph-edge count with parent interval \([0,3]\). The two
children

\[
E_{2,2}\le0
\qquad\text{and}\qquad
E_{2,2}\ge1
\tag{5.11}
\]

are disjoint and cover every possible integer value. All four child
relaxations have exact Farkas certificates. Hence the number of exact
order-20 leaves is

\[
63{,}235+16{,}542+343+4=80{,}124.
\tag{5.12}
\]

Every source at each layer is the disjoint union of the certified leaves
and the next survivor ledger. The verifier checks exact set difference,
uniqueness, source order, shard hashes and zero uncovered cases.
Therefore no retained candidate survives the certified computation.
Lemma 4 below supplies the canonical complement and graph-to-model
implications that complete the proof of Theorem 2.

\subsubsection{Computational Proposition 2 (order-20 exact
coverage)}\label{computational-proposition-2-order-20-exact-coverage}

The order-20 ledgers partition all \(80{,}122\) complement
representatives into \(63{,}235\) Stage-1 leaves, \(16{,}542\) Stage-2
leaves, \(343\) Stage-3 leaves and two final roots. The four certified
children of those roots complete the \(80{,}124\)-leaf exhaustion.

\subsection{6. Order-specific graph-to-model soundness and theorem
proofs}\label{order-specific-graph-to-model-soundness-and-theorem-proofs}

The shared identities do not make the two finite computations identical.
The order-16 proof branches through complete mixed-integer class trees,
whereas the order-20 proof first selects one complement representative,
uses continuous class relaxations, and consults integrality only for its
final split. We therefore state the two mapping lemmas separately.

\subsubsection{Lemma 3 (order-16 graph-to-model
mapping)}\label{lemma-3-order-16-graph-to-model-mapping}

Suppose that a simple graph \(G\) on \(16\) vertices has Laplacian
spectrum \(\{0,1,\ldots,15\}\). Its sorted degree sequence and forced
triangle count occur in the complete order-16 enumeration. Its
projector-diagonal weights \(Y\) satisfy the Stage-1 equations, and the
same weights together with the actual neighbor-degree sums \(S_d\)
satisfy every Stage-2 row. If the candidate reaches one of the 42 class
models, the same \(Y\) together with the actual integer degree-class
counts \(E,T,C\) gives a feasible mixed point of that model. At each
branch node, the actual integer count lies in exactly one child.

\paragraph{Proof}\label{proof}

Equations (2.1)--(2.8) and the order-16 specialization (4.1) put the
sorted degree sequence and its forced integral triangle count in the
initial list; the completeness and filtering statement is the order-16
part of Appendix A. Equations (3.1)--(3.3) follow by summing squared
eigenvector coordinates over each degree class. Equations (3.4)--(3.11),
after the stated elimination of \(H_d\), follow from the diagonal of
\(L^3\), actual neighbor sums, triangle capacities and complement
counting. For a final model, take \(E,T,C\) to be the actual edge and
graph/complement triangle counts by degree class. Every row in Appendix
B is then a capacity bound, double-counting identity, local third-moment
identity or incidence upper bound. Finally, every branch is on an
integer graph count and uses adjacent bounds \(x_j\le a\) and
\(x_j\ge a+1\), so the children are disjoint and exhaustive. Thus an
actual target graph remains feasible until it reaches one certified
leaf. \(\square\)

\subsubsection{Proof of Theorem 1}\label{proof-of-theorem-1}

Assume that a target graph of order 16 exists. Lemma 3 maps it to an
enumerated candidate and preserves a feasible point through every
applicable relaxation and integer branch. Computational Proposition 1
says that the exact coverage ledger sends every such candidate to a
verified Farkas leaf. The contradiction at that leaf rules out the
feasible point. Hence no target graph exists. \(\square\)

\subsubsection{Lemma 4 (order-20 canonical graph-to-model
mapping)}\label{lemma-4-order-20-canonical-graph-to-model-mapping}

Suppose that a simple graph \(G\) on \(20\) vertices has Laplacian
spectrum \(\{0,1,\ldots,19\}\). Let \(G^{\ast}\) be \(G\) if the pair
consisting of its sorted degree sequence and triangle count is the
retained representative of its \(\kappa\)-orbit, and let
\(G^{\ast}=\overline G\) otherwise. Then \(G^{\ast}\) realizes exactly
one of the \(80{,}122\) retained representatives. Its projector-diagonal
weights and actual neighbor-degree sums give feasible points of its
Stage-1 and Stage-2 models. If it reaches Stage 3, its actual
degree-class counts \(E,T,C\), viewed in the continuous relaxation, give
a feasible point of the class model. If it is one of the final two
roots, its integer \(E_{2,2}\) lies in exactly one of the two children
in (5.11).

\paragraph{Proof}\label{proof-1}

The target spectrum is invariant under complementation, and (5.3) is a
fixed-point-free involution of the complete \(160{,}244\)-candidate set.
Consequently exactly one of \(G,\overline G\) realizes the retained
member of its orbit; this is the complement-coverage statement
formalized in Appendix A. Applying (3.1)--(3.11) to \(G^{\ast}\) proves
the Stage-1 and Stage-2 claims exactly as above. At Stage 3 take
\(E,T,C\) to be the actual graph and complement counts by original
degree class. The complete rows in Appendix C are satisfied by the
corresponding capacity, stub, triangle-total, graph/complement
third-moment, incidence and wedge counts. Although these counts are
integers for \(G^{\ast}\), the 343 ordinary class certificates discard
their integrality and prove infeasibility of the larger continuous
relaxation. For either final root, \(E_{2,2}\) is an integer in
\([0,3]\), so it satisfies exactly one of \(E_{2,2}\le0\) and
\(E_{2,2}\ge1\). Thus \(G^{\ast}\) remains feasible until it reaches one
certified leaf. \(\square\)

\subsubsection{Proof of Theorem 2}\label{proof-of-theorem-2}

Assume that a target graph of order 20 exists and form \(G^{\ast}\) as
in Lemma 4. Computational Proposition 2 and the exact set-partition
ledger place its retained candidate in a Stage-1, Stage-2 or Stage-3
Farkas leaf, or in one of the four certified children of the final two
roots. Lemma 4 supplies a feasible point in that same leaf,
contradicting its exact certificate. Hence no target graph exists.
\(\square\)

\subsection{7. Artifacts, hashes and
reproduction}\label{artifacts-hashes-and-reproduction}

The theorem claims depend on two frozen proof artifacts. The canonical
self-contained reproduction command from the combined release root is

\begin{verbatim}
./verification/REPRODUCE_BOTH.sh
\end{verbatim}

It checks the top-level manifest, replays order 16, performs the
hardened offline order-20 replay, and requires the reproduced order-20
canonical JSON to match the frozen expected bytes.

A successful run terminates with

\begin{verbatim}
COMBINED_REPRODUCTION_PASS orders=16,20
\end{verbatim}

\subsubsection{7.1 Order-16 artifact}\label{order-16-artifact}

\begin{verbatim}
snn16-certified-nonexistence-20260821.zip
SHA-256: f7f127da4fd6227bd66eadfc22847da270caf2d29b8bedbc16aabee548c8c847
\end{verbatim}

From the extracted package root, run:

\begin{verbatim}
./REPRODUCE.sh
\end{verbatim}

The two primary certificate files have SHA-256 identities:

\begin{verbatim}
certs/degree_reduction.json.gz
950d05274bf8e7234c0c72efa78927feeeacadc33ba5274f42c738ff3e8562c9

certs/class_exhaustion.json.gz
51cc5417f922f6bfceb3704c81a1691d813f76550517397269404b35225cc9bd
\end{verbatim}

Expected statuses are VERIFIED\_UNSAT and SEMANTIC\_AUDIT\_PASSED.

\subsubsection{7.2 Order-20 artifact}\label{order-20-artifact}

\begin{verbatim}
S20_CERTIFIED_EXHAUSTION_RELEASE_20260821.tar.gz
SHA-256: 1eeda59a36dc835ec0efd3dc741d985145054af0d72e77f59e84f9cb63461206
\end{verbatim}

To replay order 20 alone from the combined release root, run:

\begin{verbatim}
./verification/n20/REPRODUCE_INDEPENDENT.sh \
  artifacts/S20_CERTIFIED_EXHAUSTION_RELEASE_20260821.tar.gz \
  tmp/reproduced-n20-canonical.json
\end{verbatim}

The publication-revision canonical JSON has SHA-256:

\begin{verbatim}
ebc3d8065c3f0c2e869e940f2d50c572292c45646a5ee6ac658a3f805b1ccbac
\end{verbatim}

Expected statuses include:

\begin{verbatim}
VERIFIED_UNSAT
INDEPENDENT_ENUMERATION_V2_PASS
INDEPENDENTLY_VERIFIED_UNSAT
GRAPH_IDENTITY_CYCLE_V2_PASS
INDEPENDENT_REPRODUCTION_PASS
\end{verbatim}

The strings beginning with INDEPENDENT are retained machine labels for
separately implemented internal checks. They do not claim external
replication.

\subsubsection{7.3 Trust boundary}\label{trust-boundary}

{\def\LTcaptype{none} % do not increment counter
\begin{longtable}[]{@{}
  >{\raggedright\arraybackslash}p{(\linewidth - 2\tabcolsep) * \real{0.5000}}
  >{\raggedright\arraybackslash}p{(\linewidth - 2\tabcolsep) * \real{0.5000}}@{}}
\toprule\noalign{}
\begin{minipage}[b]{\linewidth}\raggedright
Component
\end{minipage} & \begin{minipage}[b]{\linewidth}\raggedright
Trust role
\end{minipage} \\
\midrule\noalign{}
\endhead
\bottomrule\noalign{}
\endlastfoot
floating-point LP and generators & untrusted discovery only \\
analytic identities and mapping lemma & trusted mathematics \\
exact enumerators & trusted internal implementation \\
exact certificate verifiers & trusted internal implementation \\
coverage and branch audits & trusted internal implementation \\
external or different-language checker & not yet available; not
claimed \\
\end{longtable}
}

The two theorem packages use different constants, candidate ledgers and
final certificate structures. Passing one verifier is not treated as
evidence for the other theorem.

\subsection{8. Limitations and external
validation}\label{limitations-and-external-validation}

Theorems 1 and 2 concern only the two stated orders and do not settle
the conjecture in general. External specialist review, independently
written verifiers, clean-room reproduction on additional platforms and
eventual proof-assistant formalization would provide further assurance
and remain future work.

\newpage

\subsection{Appendix A: Exact enumeration and finite-coverage
specification}\label{appendix-a-exact-enumeration-and-finite-coverage-specification}

This appendix specifies the discrete enumeration and coverage
obligations used in the order-16 and order-20 proofs. The counts and
hashes below were checked against the frozen inputs identified in
Section A.1. All arithmetic in the trusted enumeration and coverage
checks is integer arithmetic.

\subsubsection{A.1 Frozen inputs and derived
parameters}\label{a.1-frozen-inputs-and-derived-parameters}

The two proof archives are logically separate.

\begin{verbatim}
snn16-certified-nonexistence-20260821.zip
SHA-256: f7f127da4fd6227bd66eadfc22847da270caf2d29b8bedbc16aabee548c8c847

S20_CERTIFIED_EXHAUSTION_RELEASE_20260821.tar.gz
SHA-256: 1eeda59a36dc835ec0efd3dc741d985145054af0d72e77f59e84f9cb63461206
\end{verbatim}

For a target graph of order \(n\), write its nondecreasing degree
sequence as \(s=(d_1,\ldots,d_n)\). The analytic part of the proof gives
the following enumeration parameters.

{\def\LTcaptype{none} % do not increment counter
\begin{longtable}[]{@{}lrr@{}}
\toprule\noalign{}
Parameter & \(n=16\) & \(n=20\) \\
\midrule\noalign{}
\endhead
\bottomrule\noalign{}
\endlastfoot
allowed degree interval & \(2,\ldots,13\) & \(2,\ldots,17\) \\
sequence length & 16 & 20 \\
forced degree sum & 120 & 190 \\
forced square sum & 1120 & 2280 \\
constant \(C_n\) in \(\sum_i d_i^3-6t=C_n\) & 11040 & 29260 \\
allowed triangle interval & \(54,\ldots,166\) & \(127,\ldots,348\) \\
\(t(G)+t(\overline G)\) & 220 & 475 \\
\end{longtable}
}

Consequently, if the allowed degree alphabet is
\(D_n=\{a,a+1,\ldots,b\}\), every admissible sorted sequence corresponds
bijectively to a nonnegative integer multiplicity vector
\((m_a,\ldots,m_b)\) satisfying

\[
\sum_{d=a}^{b}m_d=n,
\qquad
\sum_{d=a}^{b}d m_d=\sigma_1,
\qquad
\sum_{d=a}^{b}d^2m_d=\sigma_2.
\tag{A.1}
\]

No graph-isomorphism reduction or heuristic symmetry breaking is used in
this enumeration.

\subsubsection{A.2 Complete multiplicity
recursion}\label{a.2-complete-multiplicity-recursion}

The following pseudocode is a mathematical specification of a complete
multiplicity enumeration. The frozen order-20 verifier implements its
memoized exact-completion version. The separately structured order-16
audit implements the same multiplicity principle with endpoint pruning;
the order-16 strict verifier instead traverses every sorted tuple
directly.

\begin{verbatim}
ENUM(d, r, S, Q, multiplicities):
    # d is the smallest still available degree.
    # A completion must use r entries, with residual sum S
    # and residual square sum Q, all drawn from {d,...,b}.

    if r < 0 or S < 0 or Q < 0:
        return

    if d = b:
        if S = r*b and Q = r*b^2:
            set m_b = r
            output the sorted sequence encoded by m_a,...,m_b
        return

    if S < r*d or S > r*b:
        return
    if Q < r*d^2 or Q > r*b^2:
        return
    if r > 0 and r*Q < S^2:
        return
    if r = 0:
        output only if S = Q = 0
        return

    for c = 0,...,r:
        set m_d = c
        ENUM(d+1, r-c, S-c*d, Q-c*d^2, multiplicities)
\end{verbatim}

The initial call is

\[
\operatorname{ENUM}(a,n,\sigma_1,\sigma_2,()).
\tag{A.2}
\]

The pruning invariant is the following: at a call
\(\operatorname{ENUM}(d,r,S,Q)\), a completion, if one exists, consists
of exactly \(r\) values in \([d,b]\), with sum \(S\) and square sum
\(Q\). Therefore it necessarily satisfies

\[
rd\le S\le rb,
\qquad
rd^2\le Q\le rb^2,
\qquad
rQ\ge S^2.
\tag{A.3}
\]

The last inequality is Cauchy--Schwarz. Hence every endpoint or
Cauchy--Schwarz rejection is necessary and cannot discard a completion.
At an unpruned node, the loop tries every possible value of \(m_d\).
Induction on the number of remaining degree values proves that every
solution of (A.1) is output exactly once. The order-20 verifier
additionally memoizes the exact Boolean recurrence

\[
P(d,r,S,Q)
=\bigvee_{c=0}^{r}P(d+1,r-c,S-cd,Q-cd^2),
\tag{A.4}
\]

with the evident terminal condition. This can prune more nodes than
(A.3), but it is an exact restatement of the same exhaustive case split.

\subsubsection{A.3 Graphicality and forced triangle
filtering}\label{a.3-graphicality-and-forced-triangle-filtering}

Every sequence output by (A.2) is tested for simple-graph graphicality.
The strict paths use the Erdős--Gallai inequalities. If
\(e_1\ge\cdots\ge e_n\) is the sequence in decreasing order, the test is

\[
\sum_{i=1}^{k}e_i
\le k(k-1)+\sum_{i=k+1}^{n}\min(e_i,k)
\quad(1\le k\le n).
\tag{A.5}
\]

Only graphical sequences are retained. For each such sequence, the
triangle count required by the third spectral moment is computed, rather
than searched, as

\[
t=\frac{\sum_i d_i^3-C_n}{6}.
\tag{A.6}
\]

The numerator must be divisible by 6, and the resulting integer must lie
in the interval in Section A.1. These are necessary conditions for a
target graph. Thus the moment recursion, graphicality test, divisibility
test and triangle-range test cannot remove the degree data of a target
graph.

\subsubsection{A.4 Order-16 enumeration and
coverage}\label{a.4-order-16-enumeration-and-coverage}

The order-16 strict verifier enumerates all

\[
d_1,\ldots,d_{16}\in\{2,\ldots,13\},
\qquad d_1\le\cdots\le d_{16},
\]

by a direct combinations-with-repetition traversal. It then applies the
two moment equalities, Erdős--Gallai, and (A.6). Its exact ledger is:

{\def\LTcaptype{none} % do not increment counter
\begin{longtable}[]{@{}lrrr@{}}
\toprule\noalign{}
Order-16 layer & Source & Certified leaves & Survivors \\
\midrule\noalign{}
\endhead
\bottomrule\noalign{}
\endlastfoot
two degree moments & all sorted tuples & -- & 2,249 \\
Erdős--Gallai & 2,249 & -- & 2,233 \\
triangle integrality and range & 2,233 & -- & 2,077 \\
Stage 1 spectral model & 2,077 & 1,359 & 718 \\
Stage 2 local model & 718 & 676 & 42 \\
integer class trees & 42 roots & 108 branch-tree leaves & 0 \\
\end{longtable}
}

The file \texttt{certs/degree\_reduction.json.gz} has one record, in
exact base-list order, for each of the 2,077 candidates. Each record has
exactly one of the mutually exclusive kinds \texttt{stage1\_farkas},
\texttt{stage2\_farkas}, or \texttt{survivor}. Writing \(B_{16}\) for
the ordered base list, \(L_{16,1}\) and \(L_{16,2}\) for the two
certified subsets, and \(R_{16}\) for the final survivors, the verifier
checks

\[
B_{16}=L_{16,1}\mathbin{\dot\cup}L_{16,2}
              \mathbin{\dot\cup}R_{16},
\quad
(|L_{16,1}|,|L_{16,2}|,|R_{16}|)=(1359,676,42).
\tag{A.7}
\]

Equivalently, the Stage-1 survivor set has size \(676+42=718\), and is
the disjoint union of the Stage-2 leaves and the 42 class roots. The
class certificate contains exactly those 42 roots, in the same order.
Its integer branch trees have 174 nodes: 66 branch nodes and 108 exact
Farkas leaves, with maximum depth 10. At each branch, adjacent integer
bounds \(x_j\le q\) and \(x_j\ge q+1\) are disjoint and exhaustive.

The frozen order-16 certificate identities are:

\begin{verbatim}
certs/degree_reduction.json.gz
SHA-256: 950d05274bf8e7234c0c72efa78927feeeacadc33ba5274f42c738ff3e8562c9

certs/class_exhaustion.json.gz
SHA-256: 51cc5417f922f6bfceb3704c81a1691d813f76550517397269404b35225cc9bd
\end{verbatim}

The semantic audit follows a structurally different route: it enumerates
multiplicities by the residual recursion of Section A.2 and uses
largest-degree Havel--Hakimi instead of the strict verifier's
sorted-tuple traversal and Erdős--Gallai test. It reproduces 2,249,
2,233 and 2,077, compares the complete ordered base list with all
certificate records, and checks that the class items equal the 42
survivors. This is a separately structured internal audit, not external
replication.

\subsubsection{A.5 Order-20 complement
quotient}\label{a.5-order-20-complement-quotient}

Before certificate elimination, the order-20 computation applies a
complement quotient. For a candidate \((s,t)\), define

\[
\kappa(s,t)=
\left(\operatorname{sort}(19-d_1,\ldots,19-d_{20}),475-t\right).
\tag{A.8}
\]

This map preserves the complete candidate set. Indeed, complementing a
simple graph realizing a graphical degree sequence gives a simple graph
with degree sequence \(\operatorname{sort}(19-d_i)\). Substitution in
the first two moment equations preserves \(\sum d_i=190\) and
\(\sum d_i^2=2280\). Expanding \(\sum_i(19-d_i)^3\), and using those two
moments, gives

\[
\sum_i(19-d_i)^3
=20\cdot19^3-3\cdot19^2\cdot190+3\cdot19\cdot2280-\sum_i d_i^3
=61370-\sum_i d_i^3.
\]

Therefore (A.6) gives

\[
\overline t
=\frac{61370-\sum_i d_i^3-29260}{6}
=475-\frac{\sum_i d_i^3-29260}{6}
=475-t.
\]

The interval \([127,348]\) is invariant under \(t\mapsto475-t\).
Therefore the complement mate passes the same moment, graphicality,
divisibility and range predicates.

Applying (A.8) twice returns the original sorted pair, so \(\kappa\) is
an involution. It has no fixed point: a fixed pair would require
\(t=475-t\), or \(2t=475\), which is impossible for integer \(t\). Hence
all 160,244 candidates lie in two-element orbits. Retaining the
lexicographically smaller member of every orbit leaves exactly 80,122
representatives and loses no hypothetical target graph, because the
graph or its complement realizes the retained member.

\subsubsection{A.6 Order-20 enumeration, hashes and
coverage}\label{a.6-order-20-enumeration-hashes-and-coverage}

The order-20 strict route uses the exact memoized multiplicity
recurrence (A.4), Erdős--Gallai, (A.6), and the complement quotient
(A.8). Its ledger is:

{\def\LTcaptype{none} % do not increment counter
\begin{longtable}[]{@{}lrrr@{}}
\toprule\noalign{}
Order-20 layer & Source & Certified leaves & Survivors \\
\midrule\noalign{}
\endhead
\bottomrule\noalign{}
\endlastfoot
two degree moments & multiplicity recursion & -- & 209,932 \\
Erdős--Gallai & 209,932 & -- & 200,108 \\
triangle integrality and range & 200,108 & -- & 160,244 \\
complement quotient & 160,244 & symmetry pairing & 80,122 \\
Stage 1 spectral model & 80,122 & 63,235 & 16,887 \\
Stage 2 local model & 16,887 & 16,542 & 345 \\
Stage 3 class model & 345 & 343 & 2 roots \\
final integer split & 2 roots & 4 child leaves & 0 \\
\end{longtable}
}

The following SHA-256 values hash the complete \textbf{uncompressed
canonical ledger bytes} named immediately above them. They are not
hashes of individual certificates.

\begin{verbatim}
moment_sequences.txt (209,932)
c60854d47d7c6492577999103aba99f00e50d076f38c3f4d3ad486e25c557856

graphical_sequences.txt (200,108)
7b7cc252551d8e0afdf76e9b47df007ac292b192c3474031201f4026936ef731

base_sequences.txt (160,244)
8523a5c0c0bc1c4f2d6e76be0979f984954b393393e483f4047529b0d48245e2

complement_representatives.txt / Stage-1 source (80,122)
851b058e4c5f933e22bbd61051293342e5788fd3a524fa295266dcde4633e0d4

Stage-1 survivors / Stage-2 source (16,887)
a325883c2b1d7919f70ac82d2460733f1fe9b3e373f633b42a1b6ae512fc3ad4

Stage-2 survivors / class source (345)
dd39bed2528cea440158a0161451dc29dbbb2b36fde0f91ae6558ee7424d8798

class survivors / final roots (2)
b3b6352252492f17517e00f25307f8370daec1eb50786d598218c8f540859ecc
\end{verbatim}

For each certificate stage \(i\), let \(S_i\) be its ordered source
ledger, \(R_i\) its ordered survivor ledger, and \(L_i\) the source
indices carrying certificates. The verifier checks more than the counts:
every certificate index is in range and globally strictly increasing,
every survivor is unique, the survivor list is the source-induced order,
and

\[
S_i=L_i\mathbin{\dot\cup}R_i,
\qquad
L_i=S_i\setminus R_i.
\tag{A.9}
\]

It also compares the bytes of each next-stage source with the preceding
survivor ledger:

\[
S_1=\text{complement representatives},
\quad S_2=R_1,
\quad S_3=R_2.
\tag{A.10}
\]

The two Stage-3 survivors are byte-identical to the two roots in the
final branch metadata. Each root is split on the actual integer graph
count \(E_{2,2}\) into \(E_{2,2}\le0\) and \(E_{2,2}\ge1\). These
children are disjoint and exhaustive, and all four child models have
exact certificates. Thus the certificate-leaf total is

\[
63235+16542+343+4=80124,
\]

with zero uncovered candidates.

\subsubsection{A.7 Structurally different internal audit
paths}\label{a.7-structurally-different-internal-audit-paths}

The agreement checks deliberately change both enumeration and
graphicality algorithms.

{\def\LTcaptype{none} % do not increment counter
\begin{longtable}[]{@{}
  >{\raggedright\arraybackslash}p{(\linewidth - 4\tabcolsep) * \real{0.3333}}
  >{\raggedright\arraybackslash}p{(\linewidth - 4\tabcolsep) * \real{0.3333}}
  >{\raggedright\arraybackslash}p{(\linewidth - 4\tabcolsep) * \real{0.3333}}@{}}
\toprule\noalign{}
\begin{minipage}[b]{\linewidth}\raggedright
Order
\end{minipage} & \begin{minipage}[b]{\linewidth}\raggedright
Main exact path
\end{minipage} & \begin{minipage}[b]{\linewidth}\raggedright
Structurally different internal path
\end{minipage} \\
\midrule\noalign{}
\endhead
\bottomrule\noalign{}
\endlastfoot
16 & sorted tuples by combinations with repetition; Erdős--Gallai &
degree-multiplicity residual recursion; largest-degree Havel--Hakimi \\
20, bundled & degree-multiplicity exact-completion recursion;
Erdős--Gallai & nondecreasing position-by-position feasibility dynamic
program; largest-degree Havel--Hakimi \\
20, second verifier & does not import or execute release code &
degree-alphabet meet-in-the-middle hash join; smallest-degree
Kleitman--Wang layoff \\
\end{longtable}
}

The order-20 meet-in-the-middle audit splits the degree alphabet,
indexes high-half multisets by \texttt{(length,\ sum,\ square\ sum)},
joins complementary low-half keys, and reconstructs the complete
canonical bytes. It therefore does not reuse the multiplicity recursion.
Its smallest-degree Kleitman--Wang layoff also differs from both
Erdős--Gallai and the bundled largest-degree Havel--Hakimi path.

These checks reduce the risk of a shared algorithmic mistake, but they
remain internal Python implementations using the same frozen
mathematical specification and data. Section 7.3 states the
corresponding trust boundary.

\subsection{Appendix B: Complete computational specification for the
order-16
proof}\label{appendix-b-complete-computational-specification-for-the-order-16-proof}

This appendix specifies the finite computation contained in the frozen
order-16 archive \texttt{snn16-certified-nonexistence-20260821.zip}. The
mathematical implication used throughout is

\[
\text{target graph}
\Longrightarrow
\text{feasible point of a necessary relaxation}.
\]

Consequently, exact infeasibility of a relaxation rules out the graph.
The converse implication is never assumed.

\subsubsection{B.1 Frozen identities and candidate
enumeration}\label{b.1-frozen-identities-and-candidate-enumeration}

Let \(G\) be a simple graph on 16 vertices with hypothetical Laplacian
spectrum \(\{0,1,\ldots,15\}\). The spectral trace identities, the local
degree bound, and the graph/complement triangle identity give

\[
\sum_v d_v=120,\qquad
\sum_v d_v^2=1120,\qquad
\sum_v d_v^3-6t=11040,
\tag{B.1}
\]

\[
2\le d_v\le13,\qquad
54\le t\le166,\qquad
t(G)+t(\overline G)=220.
\tag{B.2}
\]

The strict verifier enumerates every nondecreasing tuple in the
lexicographic order induced by the Python standard-library
combinations-with-replacement iterator over degrees \(2,\ldots,13\),
\(s=(d_1,\ldots,d_{16})\) satisfying the first two equations in (B.1).
It then applies the Erdős--Gallai inequalities, and finally retains a
tuple precisely when

\[
t=\frac{\sum_i d_i^3-11040}{6}
\tag{B.3}
\]

is an integer in \([54,166]\). The exact ledger is

{\def\LTcaptype{none} % do not increment counter
\begin{longtable}[]{@{}lr@{}}
\toprule\noalign{}
Enumeration layer & Exact count \\
\midrule\noalign{}
\endhead
\bottomrule\noalign{}
\endlastfoot
moment tuples & \(2{,}249\) \\
Erdős--Gallai graphical tuples & \(2{,}233\) \\
graphical tuples with admissible integral \(t\) & \(2{,}077\) \\
\end{longtable}
}

The certificate records occur in exactly this final enumeration order. A
separate internal semantic audit enumerates the moment solutions by
degree multiplicities and tests graphicality by Havel--Hakimi. It
reproduces all three counts and the complete ordered list of 2,077 pairs
\((s,t)\). This is an internal cross-check using a different enumeration
organization, not an external replication.

For the remainder fix one candidate \((s,t)\). Let

\[
m_d=|\{i:d_i=d\}|,\qquad
D=\{d:m_d>0\},
\tag{B.4}
\]

with \(D\) always traversed in increasing order.

\subsubsection{B.2 Stage 1: spectral-weight
relaxation}\label{b.2-stage-1-spectral-weight-relaxation}

Choose an orthonormal Laplacian eigenbasis \(u_0,\ldots,u_{15}\), where
\(u_k\) has eigenvalue \(k\), and define

\[
Y_{d,k}=16\sum_{v:d_v=d}u_k(v)^2,
\qquad d\in D,\quad 1\le k\le15.
\tag{B.5}
\]

The \(Y_{d,k}\) are nonnegative \textbf{real} variables. No integrality,
orthostochasticity, or integral-eigenvector condition is imposed. Every
target graph satisfies

\[
\sum_{k=1}^{15}Y_{d,k}=15m_d,
\tag{B.6}
\]

\[
\sum_{k=1}^{15}kY_{d,k}=16m_dd,
\tag{B.7}
\]

\[
\sum_{k=1}^{15}k^2Y_{d,k}=16m_d(d^2+d)
\tag{B.8}
\]

for every \(d\in D\), and

\[
\sum_{d\in D}Y_{d,k}=16
\tag{B.9}
\]

for every \(1\le k\le15\). Stage 1 consists exactly of (B.6)--(B.9) and
\(Y_{d,k}\ge0\). Exact rational Farkas certificates eliminate 1,359 of
the 2,077 candidates, leaving 718.

\subsubsection{B.3 Stage 2: local-third-moment
relaxation}\label{b.3-stage-2-local-third-moment-relaxation}

Stage 2 retains all Stage-1 variables and equations and adds one
nonnegative real variable \(S_d\) for each \(d\in D\). Semantically,

\[
S_d=\sum_{v:d_v=d}\sum_{u\sim v}d_u.
\tag{B.10}
\]

For one occurrence of \(d\), delete that occurrence from \(s\). Let
\(\ell_d\) and \(u_d\) be the sums of the \(d\) smallest and \(d\)
largest entries of the remaining 15-tuple. Then

\[
m_d\ell_d\le S_d\le m_du_d.
\tag{B.11}
\]

Put

\[
b_d=d^3+2d^2,\qquad
K_d=\sum_{k=1}^{15}k^3Y_{d,k}.
\tag{B.12}
\]

There is no separately serialized \(H_d\) variable. It is the derived
real quantity

\[
H_d=\frac{K_d}{16}-m_db_d
=\sum_{v:d_v=d}(s_v-2q_v).
\tag{B.13}
\]

The exact Stage-2 inequalities are

\[
m_dd\le H_d\le m_dd(16-d),
\tag{B.14}
\]

\[
H_d\le S_d\le H_d+2m_d\binom d2,
\tag{B.15}
\]

and, from the local complement-triangle identity,

\[
m_d\left(120-2\binom{15-d}{2}\right)
\le S_d+H_d\le120m_d.
\tag{B.16}
\]

The final Stage-2 equality is

\[
\sum_{d\in D}S_d=1120.
\tag{B.17}
\]

For avoidance of sign or scaling ambiguity, the verifier serializes the
six base inequality rows for each increasing \(d\) in this order:

\[
-K_d\le-16m_d(b_d+d),
\tag{B.18}
\]

\[
K_d\le16m_d\{b_d+d(16-d)\},
\tag{B.19}
\]

\[
K_d-16S_d\le16m_db_d,
\tag{B.20}
\]

\[
-K_d+16S_d\le32m_d\binom d2-16m_db_d,
\tag{B.21}
\]

\[
-K_d-16S_d
\le-16m_d\left\{b_d+120-2\binom{15-d}{2}\right\},
\tag{B.22}
\]

\[
K_d+16S_d\le16m_d(b_d+120).
\tag{B.23}
\]

Variable-bound rows are appended later as specified in Section B.7.
Exact rational Farkas certificates eliminate 676 of the 718 Stage-1
survivors. The remaining 42 candidates are the ordered records in
data/survivors42\_exact.json.

\subsubsection{B.4 Final degree-class
model}\label{b.4-final-degree-class-model}

For \(a,b\in D\), \(a\le b\), define the available unordered vertex-pair
capacity

\[
P_{a,b}=
\begin{cases}
\binom{m_a}{2},&a=b,\\
m_am_b,&a<b.
\end{cases}
\tag{B.24}
\]

For a nondecreasing degree triple \(z=(a,b,c)\), let \(\mu_d(z)\) be the
multiplicity of \(d\) in \(z\) and define

\[
P_z=\prod_{d\in D}\binom{m_d}{\mu_d(z)}.
\tag{B.25}
\]

Only triples with \(P_z>0\) are included. The model contains:

\begin{itemize}
\tightlist
\item
  the nonnegative real variables \(Y_{d,k}\) from Stage 1;
\item
  \(E_{a,b}\), semantically the integer number of graph edges joining
  degree classes \(a,b\);
\item
  \(T_z\), semantically the integer number of graph triangles with
  original degree multiset \(z\);
\item
  \(C_z\), semantically the integer number of complement triangles whose
  vertices have original graph-degree multiset \(z\).
\end{itemize}

Thus \(Y\) is real, whereas actual graph values of \(E,T,C\) are
integers. The model declares every \(E,T,C\) variable integer and
imposes

\[
0\le E_{a,b}\le P_{a,b},\qquad
0\le T_z\le P_z,\qquad
0\le C_z\le P_z.
\tag{B.26}
\]

\paragraph{B.4.1 Stub and total-triangle
equations}\label{b.4.1-stub-and-total-triangle-equations}

For every \(d\in D\),

\[
2E_{d,d}+\sum_{b\in D\setminus\{d\}}
E_{\min(d,b),\max(d,b)}=m_dd.
\tag{B.27}
\]

The graph and complement totals are

\[
\sum_zT_z=t,\qquad
\sum_zC_z=220-t.
\tag{B.28}
\]

\paragraph{B.4.2 Graph classwise third
moments}\label{b.4.2-graph-classwise-third-moments}

Define

\[
S_d^G=2dE_{d,d}
+\sum_{b\in D\setminus\{d\}}
bE_{\min(d,b),\max(d,b)},
\tag{B.29}
\]

\[
R_d^G=\sum_z\mu_d(z)T_z.
\tag{B.30}
\]

For each \(d\in D\), the graph classwise \(L^3\) equation is

\[
\sum_{k=1}^{15}k^3Y_{d,k}
=16\left\{m_d(d^3+2d^2)+S_d^G-2R_d^G\right\}.
\tag{B.31}
\]

\paragraph{B.4.3 Complement classwise third
moments}\label{b.4.3-complement-classwise-third-moments}

An original degree-\(d\) vertex has complement degree \(15-d\), and

\[
\overline Y_{15-d,k}=Y_{d,16-k}.
\tag{B.32}
\]

There are \(P_{a,b}-E_{a,b}\) complement edges between original degree
classes \(a,b\). Therefore define

\[
S_d^{\overline G}
=2(15-d)(P_{d,d}-E_{d,d})
+\sum_{b\in D\setminus\{d\}}
(15-b)\left(
P_{\min(d,b),\max(d,b)}
-E_{\min(d,b),\max(d,b)}
\right),
\tag{B.33}
\]

\[
R_d^{\overline G}=\sum_z\mu_d(z)C_z.
\tag{B.34}
\]

The complement classwise \(L^3\) equation is

\[
\sum_{k=1}^{15}k^3Y_{d,16-k}
=16\left\{
m_d\bigl((15-d)^3+2(15-d)^2\bigr)
+S_d^{\overline G}-2R_d^{\overline G}
\right\}.
\tag{B.35}
\]

\paragraph{B.4.4 Pair-incidence
inequalities}\label{b.4.4-pair-incidence-inequalities}

For \(a\le b\), set

\[
\kappa_{a,b}(z)=
\begin{cases}
\binom{\mu_a(z)}2,&a=b,\\
\mu_a(z)\mu_b(z),&a<b.
\end{cases}
\tag{B.36}
\]

This is the number of degree-\((a,b)\) vertex pairs in a triple of type
\(z\). Every graph edge of endpoint degrees \(a,b\) lies in at most
\(\min(a,b)-1\) graph triangles. Every complement edge between the same
original classes has complement endpoint degrees \(15-a,15-b\), and
hence lies in at most \(14-\max(a,b)\) complement triangles. The exact
inequalities are

\[
\sum_z\kappa_{a,b}(z)T_z
\le(\min(a,b)-1)E_{a,b},
\tag{B.37}
\]

\[
\sum_z\kappa_{a,b}(z)C_z
\le(14-\max(a,b))(P_{a,b}-E_{a,b}).
\tag{B.38}
\]

\paragraph{B.4.5 Degree-class wedge
inequalities}\label{b.4.5-degree-class-wedge-inequalities}

Counting triangle incidences at vertices in one degree class gives

\[
\sum_z\mu_d(z)T_z\le m_d\binom d2,
\tag{B.39}
\]

\[
\sum_z\mu_d(z)C_z\le m_d\binom{15-d}{2}.
\tag{B.40}
\]

Equations (B.6)--(B.9), (B.27)--(B.28), (B.31), and (B.35), together
with (B.26) and (B.37)--(B.40), are the complete final class model. No
graph isomorphism quotient or symmetry-breaking constraint is present.

\subsubsection{B.5 Why every target graph maps into the
models}\label{b.5-why-every-target-graph-maps-into-the-models}

For a target graph, (B.5) supplies nonnegative real \(Y\) and the
spectral projector identities give (B.6)--(B.9). Actual neighbor-degree
sums supply \(S_d\); the diagonal identity

\[
(L^3)_{vv}=d_v^3+2d_v^2+s_v-2q_v
\tag{B.41}
\]

and \(s_v-2q_v=d_v+e(N(v),R_v)\) give (B.13)--(B.15). The local
complement-triangle identity gives (B.16). Actual degree-class edge and
triangle counts supply integer \(E,T,C\). Their pair and triple
capacities, stub counts, triangle totals, classwise third moments,
edge-triangle incidences, and wedge counts give every row in Section
B.4. Complement equations follow from

\[
L(\overline G)=16I-J-L(G).
\tag{B.42}
\]

Thus every target graph gives a feasible mixed point. The relaxations
may also contain points that do not come from graphs; that enlargement
is harmless for an infeasibility proof.

\subsubsection{B.6 Integer branching and leaf
semantics}\label{b.6-integer-branching-and-leaf-semantics}

At the root of each of the 42 final models, \(Y\) is real and \(E,T,C\)
have the integer semantics and bounds in (B.26). A branch node stores a
variable index \(j\) belonging to the declared integer-variable set and
an integer \(a\). Its children impose

\[
x_j\le a
\qquad\text{and}\qquad
x_j\ge a+1.
\tag{B.43}
\]

These children are disjoint and cover every integer value allowed at the
parent. The verifier propagates current lower and upper bounds
recursively and, for current bounds \(\ell_j,u_j\), requires
\(\ell_j\le a<u_j\).

In the frozen forest all 66 branch nodes happen to branch on
\(E\)-variables. This does not assume that \(T\) or \(C\) is continuous
in an actual graph. At a leaf the stored Farkas certificate proves
infeasibility of the \textbf{continuous} linear relaxation under the
accumulated branch bounds. Hence it also excludes the smaller set in
which all \(E,T,C\) are integral; no further branch is needed.

The 42 rooted trees contain

{\def\LTcaptype{none} % do not increment counter
\begin{longtable}[]{@{}lr@{}}
\toprule\noalign{}
Quantity & Exact value \\
\midrule\noalign{}
\endhead
\bottomrule\noalign{}
\endlastfoot
roots / final class models & \(42\) \\
all tree nodes & \(174\) \\
branch nodes & \(66\) \\
Farkas leaves & \(108\) \\
maximum depth & \(10\) \\
\end{longtable}
}

The full-binary-forest identity \(108=66+42\) is an additional
elementary check on the tree ledger.

\subsubsection{B.7 Farkas certificates and deterministic
serialization}\label{b.7-farkas-certificates-and-deterministic-serialization}

Every model is represented as

\[
Ax\le b,\qquad Cx=e.
\tag{B.44}
\]

Finite lower bounds are appended as rows \(-x_j\le-\ell_j\); finite
upper bounds are appended as \(x_j\le u_j\). A certificate consists of
rational multipliers \(y,z\) such that

\[
y\ge0,\qquad
y^{\mathsf T}A+z^{\mathsf T}C=0,\qquad
y^{\mathsf T}b+z^{\mathsf T}e<0.
\tag{B.45}
\]

The frozen 2,143 certificates are normalized so that the last quantity
is exactly \(-1\). The verifier checks cancellation and negativity over
fractions.Fraction, without a floating tolerance. Equation (B.45)
contradicts feasibility because a feasible \(x\) would give \(0\le-1\).

The deterministic variable order is:

\begin{enumerate}
\def\labelenumi{\arabic{enumi}.}
\tightlist
\item
  \(Y_{d,k}\), with \(d\) increasing and then \(k=1,\ldots,15\);
\item
  in Stage 2 only, \(S_d\) with \(d\) increasing;
\item
  in a class model, \(E_{a,b}\) for lexicographically ordered pairs
  \(a\le b\);
\item
  \(T_z\) for lexicographically ordered nondecreasing triples with
  \(P_z>0\);
\item
  \(C_z\) in the same triple order.
\end{enumerate}

The deterministic equality-row order is:

\begin{enumerate}
\def\labelenumi{\arabic{enumi}.}
\tightlist
\item
  for each increasing \(d\), the three rows (B.6)--(B.8);
\item
  for \(k=1,\ldots,15\), row (B.9);
\item
  in Stage 2, the single final row (B.17);
\item
  in a class model instead: stub rows (B.27) by increasing \(d\), the
  graph and complement totals (B.28), graph \(L^3\) rows (B.31) by
  increasing \(d\), and complement \(L^3\) rows (B.35) by increasing
  original \(d\).
\end{enumerate}

The deterministic base-inequality order is:

\begin{enumerate}
\def\labelenumi{\arabic{enumi}.}
\tightlist
\item
  in Stage 2, rows (B.18)--(B.23) for each increasing \(d\);
\item
  in a class model, for each lexicographic pair \(a\le b\), graph row
  (B.37) followed by complement row (B.38);
\item
  for each increasing \(d\), graph row (B.39) followed by complement row
  (B.40).
\end{enumerate}

After those base inequalities, every finite lower bound is appended in
variable index order, followed by every finite upper bound in variable
index order. This order determines the row indices referenced by a
Farkas certificate.

\paragraph{B.7.1 Degree-reduction
JSON}\label{b.7.1-degree-reduction-json}

certs/degree\_reduction.json.gz has top-level format
snn16-degree-reduction-v1 and fields

\begin{verbatim}
base_count, counts, format, records
\end{verbatim}

There are exactly 2,077 ordered records. Each record contains sequence,
triangles, and one of the kinds stage1\_farkas, stage2\_farkas, or
survivor. The first two kinds additionally contain cert; a survivor
contains no certificate. The exact kind counts are

\begin{verbatim}
stage1_farkas  1359
stage2_farkas   676
survivor         42
\end{verbatim}

\paragraph{B.7.2 Certificate JSON}\label{b.7.2-certificate-json}

A certificate is an object with arrays y and z. A nonzero rational
multiplier is encoded as

\begin{verbatim}
[row_index, numerator, denominator]
\end{verbatim}

with nonzero denominator. Omitted row indices mean multiplier zero. The
y array indexes the complete inequality list, including bounds, and must
be nonnegative. The z array indexes the equality list and is
unrestricted in sign. Duplicate or out-of-range indices are rejected by
the strict verifier.

\paragraph{B.7.3 Class-exhaustion
JSON}\label{b.7.3-class-exhaustion-json}

certs/class\_exhaustion.json.gz has top-level format
snn16-class-exhaustion-v1, count 42, and an ordered items array. Every
item records its index, sequence, triangle count, statistics and root
tree. A tree node has one of the forms

\begin{verbatim}
{"type":"leaf","cert":{"y":[],"z":[]}}
\end{verbatim}

or

\begin{verbatim}
{
  "type":"branch",
  "var":143,
  "floor":4,
  "left":{},
  "right":{}
}
\end{verbatim}

Both displays are schema illustrations: actual leaf multiplier arrays
are populated, and the empty child objects are placeholders for
recursively encoded nodes. The verifier reconstructs the model from the
sequence and triangle count; no matrix stored in the certificate is
trusted.

\subsubsection{B.8 Exact coverage and frozen
identities}\label{b.8-exact-coverage-and-frozen-identities}

The degree-reduction records and class forest give the disjoint coverage
ledger

\[
2077=1359+676+42,
\tag{B.46}
\]

\[
2143=1359+676+108.
\tag{B.47}
\]

The strict verifier matches every degree record to the freshly
enumerated \((s,t)\) at the same position, matches the 42 class items to
exactly the 42 survivors, checks every branch recursively, and verifies
all 2,143 rational contradictions. The semantic audit independently
reorganizes the finite degree-sequence enumeration, checks the survivor
and complement ledgers, tests selected non-spectral graph-to-model
identities on structured and deterministic random graphs, and confirms
that a one-unit certificate mutation is rejected.

The frozen artifact identities are:

\begin{verbatim}
snn16-certified-nonexistence-20260821.zip
f7f127da4fd6227bd66eadfc22847da270caf2d29b8bedbc16aabee548c8c847

certs/degree_reduction.json.gz
950d05274bf8e7234c0c72efa78927feeeacadc33ba5274f42c738ff3e8562c9

certs/class_exhaustion.json.gz
51cc5417f922f6bfceb3704c81a1691d813f76550517397269404b35225cc9bd

data/survivors42_exact.json
ca899c72bf4e17c18663ee5b0d58a044347dcaad078c1a37f02bf551c1a09fa7

code/verify_certificates.py
7fdc4932f8f6a50b7bef06db8537c3eb621ae5cea9ad6910c53ea455a0458d61

code/audit_semantics.py
a6908f2f844c35cc806e356cef121ff2984662c022905276a6b910d7dcaf026c
\end{verbatim}

A clean replay from the extracted ZIP is

\begin{verbatim}
./REPRODUCE.sh
\end{verbatim}

and has expected terminal statuses VERIFIED\_UNSAT and
SEMANTIC\_AUDIT\_PASSED. The strict verifier uses Python
standard-library arithmetic only. The archive records successful normal
and optimized Python verification; certificate regeneration is not part
of the trusted replay.

The combined release additionally stores a current normal/optimized
transcript:

\begin{verbatim}
logs/n16-current.log
\end{verbatim}

The combined replay normalizes away elapsed time, requires the two
verifier records to agree exactly, and requires the resulting summary
bytes to equal:

\begin{verbatim}
results/ORDER16_RESULT_FINAL.json
\end{verbatim}

\subsubsection{B.9 Generator/verifier lineage and limits of
independence}\label{b.9-generatorverifier-lineage-and-limits-of-independence}

The untrusted generator uses code/exact\_models.py, SciPy/HiGHS and
SymPy. Floating-point optimization is used only to discover supports and
branch choices; proposed multipliers are exactified before storage. None
of the solver success or infeasibility flags is a premise of the
theorem.

The trusted strict verifier is a separately implemented
\textbf{internal} Python program. It uses only the standard library,
re-enumerates the candidate list, reconstructs every row from the
mathematical definitions, and checks the stored rational multipliers and
integer branches. It does not import the generator's exact\_models.py or
farkas\_tools.py.

The semantic audit supplies another enumeration organization,
Havel--Hakimi graphicality, selected graph-identity tests, coverage
comparisons and mutation rejection. Its mutation test loads the strict
verifier. Both programs are internal Python implementations; Section 7.3
states their trust boundary.

The remaining trusted boundary consists of the analytic graph-to-model
mapping in Section B.5, the correctness of the internal Python verifier
and runtime, and integrity of the frozen hashes. Those limitations
affect the validation status of the computer-assisted proof but do not
turn floating-point discovery output into a trusted premise.

\subsection{Appendix C: Complete order-20 model and certificate
specification}\label{appendix-c-complete-order-20-model-and-certificate-specification}

\subsubsection{C.1 Purpose, scope and frozen
artifact}\label{c.1-purpose-scope-and-frozen-artifact}

This appendix gives the complete mathematical and serialization
specification for the order-\(20\) theorem in the combined manuscript.
It makes the graph-to-model implication and the correspondence with the
exact verifier auditable without treating generated matrices or
floating-point solver output as trusted evidence.

The theorem addressed here is

\[
\text{there is no simple graph }G\text{ on }20\text{ vertices with }
\operatorname{spec}L(G)=\{0,1,\ldots,19\}.
\]

The frozen certificate artifact is

\begin{verbatim}
S20_CERTIFIED_EXHAUSTION_RELEASE_20260821.tar.gz
SHA-256: 1eeda59a36dc835ec0efd3dc741d985145054af0d72e77f59e84f9cb63461206
\end{verbatim}

Throughout, \(G\) is a finite simple undirected graph, \(A\) is its
adjacency matrix, \(D=\operatorname{diag}(d_v)\), and \(L=D-A\). Every
model below contains necessary conditions only. A feasible relaxation
point need not come from a graph; soundness requires the other
direction: every hypothetical target graph supplies a feasible point.

\subsubsection{C.2 Global necessities, enumeration and complement
quotient}\label{c.2-global-necessities-enumeration-and-complement-quotient}

The prescribed spectrum has a simple zero, so \(G\) is connected. Its
first three Laplacian moments give

\[
\sum_vd_v=190,\qquad
\sum_vd_v^2=2280,\qquad
\sum_vd_v^3-6t(G)=29260.
\tag{C.1}
\]

Thus \(|E(G)|=95\), and the only possible triangle count for a degree
sequence is

\[
t(G)=\frac{\sum_vd_v^3-29260}{6}.
\tag{C.2}
\]

For an orthonormal eigenbasis \(Lu_k=ku_k\), put \(w_{v,k}=u_k(v)^2\).
Since \(w_{v,0}=1/20\),

\[
\sum_{k=1}^{19}kw_{v,k}=d_v,\qquad
\sum_{k=1}^{19}k^2w_{v,k}=d_v^2+d_v,\qquad
\sum_{k=1}^{19}w_{v,k}=\frac{19}{20}.
\]

The pointwise inequality \(k^2\le20k-19\), valid for \(1\le k\le19\),
implies

\[
20d_v^2-380d_v+361\le0.
\]

Its integral solutions in the simple-graph range are exactly

\[
2\le d_v\le17.
\tag{C.3}
\]

For an edge \(uv\), the number \(c_{uv}\) of common neighbors satisfies
\(c_{uv}\ge d_u+d_v-20\). Hence

\[
3t(G)=\sum_{uv\in E(G)}c_{uv}
\ge\sum_{uv\in E(G)}(d_u+d_v-20)
=2280-20(95)=380,
\]

so \(t(G)\ge127\). The identity

\[
L(\overline G)=20I-J-L(G)
\]

shows that the target spectrum is complement-invariant. Also,

\[
t(G)+t(\overline G)
=\binom{20}{3}-\frac12\sum_vd_v(19-d_v)=475.
\tag{C.4}
\]

Applying the same lower bound to \(\overline G\) yields

\[
127\le t(G)\le348.
\tag{C.5}
\]

Exact enumeration lists every nondecreasing \(20\)-tuple in \([2,17]\)
satisfying the degree moments, tests graphicality, and imposes the
integrality and range in (C.2) and (C.5). The counts are

\[
209{,}932\longrightarrow200{,}108\longrightarrow160{,}244
\tag{C.6}
\]

for moment sequences, graphical sequences, and graphical
degree-sequence/triangle-count pairs.

For \(s=(d_1,\ldots,d_{20})\) nondecreasing, define

\[
\kappa(s,t)=
\left(\operatorname{sort}(19-d_1,\ldots,19-d_{20}),475-t\right).
\tag{C.7}
\]

The verifier checks that \(\kappa\) is an involution on the complete
\(160{,}244\)-candidate set. It has no fixed pair because a fixed pair
would require the integer equation \(2t=475\). Retaining the
lexicographically smaller member of each orbit gives \(80{,}122\)
representatives. If a target graph realizes a discarded member, its
complement realizes the retained mate.

\subsubsection{C.3 Stage 1: degree-class spectral
weights}\label{c.3-stage-1-degree-class-spectral-weights}

Let \(m_d=|\{v:d_v=d\}|\), let \(D_s\) be the increasing set of degrees
in a candidate \(s\), and define nonnegative real variables

\[
Y_{d,k}=20\sum_{v:d_v=d}u_k(v)^2,
\qquad d\in D_s,\quad1\le k\le19.
\tag{C.8}
\]

Every target graph satisfies, for each \(d\in D_s\),

\[
\sum_{k=1}^{19}Y_{d,k}=19m_d,
\tag{C.9}
\]

\[
\sum_{k=1}^{19}kY_{d,k}=20m_dd,
\tag{C.10}
\]

\[
\sum_{k=1}^{19}k^2Y_{d,k}=20m_d(d^2+d),
\tag{C.11}
\]

and, for every \(1\le k\le19\),

\[
\sum_{d\in D_s}Y_{d,k}=20.
\tag{C.12}
\]

Equations (C.9)--(C.12), together with \(Y_{d,k}\ge0\), are exactly the
Stage-1 model. Orthostochasticity, eigenvector signs, and other
nonlinear conditions are omitted, so this is a relaxation. Exact Farkas
certificates eliminate \(63{,}235\) of the \(80{,}122\) representatives
and leave \(16{,}887\).

\subsubsection{C.4 Stage 2: local third moments and all serialized
rows}\label{c.4-stage-2-local-third-moments-and-all-serialized-rows}

For a vertex \(v\) of degree \(d\), define

\[
s_v=\sum_{u\sim v}d_u,\qquad
q_v=\#\{\text{triangles of }G\text{ containing }v\},\qquad
h_v=s_v-2q_v.
\tag{C.13}
\]

Direct multiplication gives

\[
(L^3)_{vv}=d^3+2d^2+s_v-2q_v=d^3+2d^2+h_v.
\tag{C.14}
\]

If \(R_v=V\setminus(N(v)\cup\{v\})\), then

\[
h_v=d+e(N(v),R_v),\qquad
d\le h_v\le d(20-d).
\tag{C.15}
\]

For a degree class put

\[
S_d=\sum_{v:d_v=d}s_v,\qquad
Q_d=\sum_{v:d_v=d}q_v,\qquad
H_d=S_d-2Q_d.
\tag{C.16}
\]

Also put

\[
A_d=\sum_{k=1}^{19}k^3Y_{d,k},\qquad
b_d=d^3+2d^2.
\tag{C.17}
\]

Summing (C.14) over the class gives the elimination identity

\[
A_d=20(m_db_d+H_d),\qquad
H_d=\frac{A_d}{20}-m_db_d.
\tag{C.18}
\]

The serialized Stage-2 model contains \(S_d\), but it does \textbf{not}
contain \(H_d\) or \(Q_d\) as variables and does not serialize (C.18) as
an extra row. Those symbols only derive the six necessary inequalities
below.

The graph-triangle bounds give

\[
m_dd\le H_d\le m_dd(20-d),\qquad
H_d\le S_d\le H_d+2m_d\binom d2.
\tag{C.19}
\]

For a degree-\(d\) vertex, the number of complement triangles through it
is

\[
\overline q_v=\binom{19-d}{2}-95+s_v-q_v.
\tag{C.20}
\]

Consequently,

\[
m_d\left(190-2\binom{19-d}{2}\right)
\le S_d+H_d\le190m_d.
\tag{C.21}
\]

Let

\[
R_d^*=190-2\binom{19-d}{2}.
\]

After substituting (C.18), the exact six serialized inequalities are

\[
\begin{aligned}
A_d &\ge20m_d(b_d+d),\\
A_d &\le20m_d\bigl(b_d+d(20-d)\bigr),\\
A_d-20S_d &\le20m_db_d,\\
-A_d+20S_d &\le40m_d\binom d2-20m_db_d,\\
-A_d-20S_d &\le-20m_d(b_d+R_d^*),\\
A_d+20S_d &\le20m_d(b_d+190).
\end{aligned}
\tag{C.22}
\]

Delete one occurrence of \(d\) from the degree multiset and sort the
remaining entries as \(r_1\le\cdots\le r_{19}\). A degree-\(d\) neighbor
set selects \(d\) distinct entries, so

\[
m_d\sum_{i=1}^{d}r_i
\le S_d\le
m_d\sum_{i=20-d}^{19}r_i.
\tag{C.23}
\]

Oriented-edge double counting gives

\[
\sum_{d\in D_s}S_d=2280.
\tag{C.24}
\]

Stage 2 consists of Stage 1, real \(S_d\) variables with bounds (C.23),
the six inequalities (C.22), and equality (C.24). It eliminates
\(16{,}542\) of the \(16{,}887\) Stage-1 survivors and leaves \(345\).

\subsubsection{C.5 Stage 3: degree-pair and degree-triple class
model}\label{c.5-stage-3-degree-pair-and-degree-triple-class-model}

Fix one of the \(345\) Stage-2 survivors. For \(a\le b\) in \(D_s\),
define

\[
P_{a,b}=
\begin{cases}
\binom{m_a}{2},&a=b,\\
m_am_b,&a<b.
\end{cases}
\tag{C.25}
\]

For a nondecreasing triple \(z=(a,b,c)\), let \(\nu_d(z)\) be the
multiplicity of \(d\) in \(z\), and define

\[
P_z=\prod_d\binom{m_d}{\nu_d(z)}.
\tag{C.26}
\]

Triples with \(P_z=0\) are omitted. For an actual graph define:

\begin{itemize}
\tightlist
\item
  \(E_{a,b}\): the number of graph edges whose endpoint-degree multiset
  is \(\{a,b\}\);
\item
  \(T_z\): the number of graph triangles whose vertex-degree multiset is
  \(z\);
\item
  \(C_z\): the number of complement triangles whose vertices have
  original-graph degree multiset \(z\).
\end{itemize}

Their elementary capacities are

\[
0\le E_{a,b}\le P_{a,b},\qquad
0\le T_z,C_z\le P_z.
\tag{C.27}
\]

\paragraph{C.5.1 Stub and triangle
totals}\label{c.5.1-stub-and-triangle-totals}

Counting degree-\(d\) edge ends gives

\[
2E_{d,d}+\sum_{b\ne d}E_{\min(d,b),\max(d,b)}=m_dd.
\tag{C.28}
\]

The graph and complement triangle totals are

\[
\sum_zT_z=t,\qquad
\sum_zC_z=475-t.
\tag{C.29}
\]

\paragraph{C.5.2 Graph classwise third
moments}\label{c.5.2-graph-classwise-third-moments}

Define

\[
S_d^G
=2dE_{d,d}
+\sum_{b\ne d}bE_{\min(d,b),\max(d,b)},
\tag{C.30}
\]

and

\[
Q_d^G=\sum_z\nu_d(z)T_z.
\tag{C.31}
\]

The graph classwise \(L^3\) equation is

\[
\sum_{k=1}^{19}k^3Y_{d,k}
=20\left[m_d(d^3+2d^2)+S_d^G-2Q_d^G\right].
\tag{C.32}
\]

The verifier serializes (C.32) after substituting (C.30)--(C.31); those
definitions are not additional rows.

\paragraph{C.5.3 Complement classwise third
moments}\label{c.5.3-complement-classwise-third-moments}

An original degree-\(d\) vertex has complement degree \(19-d\), and

\[
\overline Y_{19-d,k}=Y_{d,20-k}.
\tag{C.33}
\]

For \(a\le b\), define

\[
\gamma_d(a,b)=
\begin{cases}
2(19-d),&a=b=d,\\
19-b,&a=d<b,\\
19-a,&a<d=b,\\
0,&d\notin\{a,b\}.
\end{cases}
\tag{C.34}
\]

There are \(P_{a,b}-E_{a,b}\) complement edges of original class pair
\((a,b)\). Hence

\[
S_d^{\overline G}
=\sum_{a\le b}\gamma_d(a,b)(P_{a,b}-E_{a,b}),
\tag{C.35}
\]

and

\[
Q_d^{\overline G}=\sum_z\nu_d(z)C_z.
\tag{C.36}
\]

Put \(\overline b_d=(19-d)^3+2(19-d)^2\). Applying the local \(L^3\)
identity to the complement and moving the \(-E_{a,b}\) terms to the left
gives the exact verifier row

\[
\sum_{k=1}^{19}k^3Y_{d,20-k}
+20\sum_{a\le b}\gamma_d(a,b)E_{a,b}
+40\sum_z\nu_d(z)C_z
=20m_d\overline b_d
+20\sum_{a\le b}\gamma_d(a,b)P_{a,b}.
\tag{C.37}
\]

The \(E_{a,b}\) coefficient is positive because the complement uses
\(P_{a,b}-E_{a,b}\).

\paragraph{C.5.4 Pair-incidence and wedge
inequalities}\label{c.5.4-pair-incidence-and-wedge-inequalities}

For \(a\le b\), define

\[
\pi_{a,b}(z)=
\begin{cases}
\binom{\nu_a(z)}2,&a=b,\\
\nu_a(z)\nu_b(z),&a<b.
\end{cases}
\tag{C.38}
\]

The graph edge--triangle incidence bound is

\[
\sum_z\pi_{a,b}(z)T_z
\le(\min(a,b)-1)E_{a,b}.
\tag{C.39}
\]

An original nonedge of class \((a,b)\) is a complement edge with
endpoint degrees \(19-a,19-b\), so

\[
\sum_z\pi_{a,b}(z)C_z
\le(18-\max(a,b))(P_{a,b}-E_{a,b}).
\tag{C.40}
\]

The graph and complement wedge bounds are

\[
\sum_z\nu_d(z)T_z\le m_d\binom d2,\qquad
\sum_z\nu_d(z)C_z\le m_d\binom{19-d}{2}.
\tag{C.41}
\]

Stage 3 comprises spectral rows (C.9)--(C.12), bounds (C.27), equations
(C.28)--(C.29), the substituted classwise rows (C.32) and (C.37), and
inequalities (C.39)--(C.41).

\paragraph{C.5.5 Continuous-relaxation and integrality
semantics}\label{c.5.5-continuous-relaxation-and-integrality-semantics}

For an actual graph, \(Y\) is generally real, whereas \(E,T,C\) are
integer graph counts. In each of the \(343\) ordinary Stage-3 leaves,
the verifier allows \(E,T,C\) to be real. Each ordinary certificate
therefore proves infeasibility of a continuous relaxation containing
every actual mixed graph point. Integrality metadata is not imposed in
those root LPs. It is consulted only to validate the final branch
variable. Root certificates eliminate \(343\) of the \(345\) models and
leave two roots.

\subsubsection{C.6 Exact Farkas
convention}\label{c.6-exact-farkas-convention}

Before standardization, write a model as

\[
A_{\rm eq}x=b_{\rm eq},\qquad
A_{\rm ub}x\le b_{\rm ub},\qquad
\ell_j\le x_j\le u_j,
\tag{C.42}
\]

where every lower bound is finite. Substitute \(x=\ell+x'\), \(x'\ge0\).
For each inequality add a nonnegative slack \(s\); for each finite upper
bound add a nonnegative slack \(r\):

\[
\begin{aligned}
A_{\rm eq}x'&=b_{\rm eq}-A_{\rm eq}\ell,\\
A_{\rm ub}x'+s&=b_{\rm ub}-A_{\rm ub}\ell,\\
x'_j+r_j&=u_j-\ell_j.
\end{aligned}
\tag{C.43}
\]

This gives an integer standard form

\[
Bw=h,\qquad w\ge0.
\tag{C.44}
\]

A sparse integer row multiplier \(z\) is a contradiction certificate
when

\[
B^{\mathsf T}z\ge0,\qquad h^{\mathsf T}z<0.
\tag{C.45}
\]

A feasible \(w\) would instead give
\(h^{\mathsf T}z=w^{\mathsf T}B^{\mathsf T}z\ge0\). Equality-row
multipliers may have either sign. Every model-inequality and upper-bound
row has its own \(+1\) slack column, so its multiplier must be
nonnegative. The checker verifies lower-bound shifts, sparse indices,
integer coefficients, slack-row signs, every component of
\(B^{\mathsf T}z\), and the strict inequality \(h^{\mathsf T}z<0\). It
reconstructs the model from candidate data and trusts neither serialized
matrices nor floating-point solver verdicts.

\subsubsection{C.7 Deterministic variable, row and certificate
serialization}\label{c.7-deterministic-variable-row-and-certificate-serialization}

For every candidate, \(D_s\) is increasing. Stage-1 variables are
\(Y_{d,k}\), first by increasing \(d\), then by \(k=1,\ldots,19\). Stage
2 uses the same block and appends \(S_d\) in increasing degree order.
The class model orders variables as:

\begin{enumerate}
\def\labelenumi{\arabic{enumi}.}
\tightlist
\item
  \(Y_{d,k}\), first by increasing \(d\), then \(k\);
\item
  \(E_{a,b}\) for lexicographically ordered pairs \(a\le b\);
\item
  \(T_z\) for lexicographically ordered nondecreasing triples with
  \(P_z>0\);
\item
  \(C_z\) in the same triple order.
\end{enumerate}

Standard-form rows are ordered as:

\begin{enumerate}
\def\labelenumi{\arabic{enumi}.}
\tightlist
\item
  equality rows in construction order;
\item
  model inequalities in construction order;
\item
  finite variable upper bounds in variable order.
\end{enumerate}

For Stage 1 and Stage 2, equality order is the three class moments
(C.9)--(C.11) for increasing \(d\), the spectral-column equations (C.12)
for increasing \(k\), and, in Stage 2, (C.24). Stage-2 inequalities are
the six rows (C.22) for each increasing \(d\). Finite \(S_d\) upper
bounds from (C.23) appear later with the variable upper-bound rows;
their lower bounds are handled by the shift in (C.43).

For the class model, equality order is:

\begin{enumerate}
\def\labelenumi{\arabic{enumi}.}
\tightlist
\item
  spectral class moments for increasing \(d\);
\item
  spectral-column equations for increasing \(k\);
\item
  stub equations for increasing \(d\);
\item
  graph and complement triangle totals;
\item
  graph classwise \(L^3\) rows for increasing \(d\);
\item
  complement classwise \(L^3\) rows for increasing \(d\).
\end{enumerate}

Class inequalities are the graph and complement pair-incidence rows for
each lexicographically ordered pair, followed by graph and complement
wedge rows for increasing \(d\). Upper-bound rows then follow the
\(E,T,C\) order.

Each Stage-1, Stage-2, or class certificate shard has schema

\begin{verbatim}
{
  "format": "s20-standard-form-farkas-z-v1",
  "stage": "stage1 | stage2 | class",
  "records": [
    ["source_index", [["row_index", "integer_multiplier"]]]
  ]
}
\end{verbatim}

The quoted field values in this display denote types, not literal values
in an actual shard. Within each certificate, nonzero row indices are
strictly increasing, unique, and in range. Across consecutively numbered
shards, source indices are globally increasing and equal the exact
source-minus-survivor index list. Zero multipliers and noninteger
coefficients are rejected.

The final branch file is certs/branch\_leaves.json.gz. For each root it
binds the source sequence and triangle count to a verified variable
index and name, the integer split floor, and one sparse certificate for
each child. Its metadata must agree with
data/final\_integer\_roots.json.

\subsubsection{C.8 Set partition, ledger identities and exact
coverage}\label{c.8-set-partition-ledger-identities-and-exact-coverage}

At every certified layer,

\[
S_i=L_i\mathbin{\dot\cup}S_{i+1},\qquad
L_i=S_i\setminus S_{i+1}.
\tag{C.46}
\]

The exact coverage ledger is:

{\def\LTcaptype{none} % do not increment counter
\begin{longtable}[]{@{}lrrr@{}}
\toprule\noalign{}
Layer & Source & Exact Farkas leaves & Survivors \\
\midrule\noalign{}
\endhead
\bottomrule\noalign{}
\endlastfoot
complement representatives & \(80{,}122\) & -- & \(80{,}122\) \\
Stage 1 & \(80{,}122\) & \(63{,}235\) & \(16{,}887\) \\
Stage 2 & \(16{,}887\) & \(16{,}542\) & \(345\) \\
Stage 3 & \(345\) & \(343\) & \(2\) roots \\
integer split & \(2\) roots & \(4\) & \(0\) \\
\end{longtable}
}

The five complete uncompressed ledger-byte identities are

\begin{verbatim}
base candidates (160,244)
8523a5c0c0bc1c4f2d6e76be0979f984954b393393e483f4047529b0d48245e2

complement representatives / Stage-1 source (80,122)
851b058e4c5f933e22bbd61051293342e5788fd3a524fa295266dcde4633e0d4

Stage-1 survivors / Stage-2 source (16,887)
a325883c2b1d7919f70ac82d2460733f1fe9b3e373f633b42a1b6ae512fc3ad4

Stage-2 survivors / class source (345)
dd39bed2528cea440158a0161451dc29dbbb2b36fde0f91ae6558ee7424d8798

class survivors / final roots (2)
b3b6352252492f17517e00f25307f8370daec1eb50786d598218c8f540859ecc
\end{verbatim}

The checker verifies ordered subset membership, uniqueness, exact set
difference, source order, shard numbering, and shard hashes. It records
zero uncovered cases, unknown outcomes, timeouts, or crashes.

\subsubsection{\texorpdfstring{C.9 Final \(E_{2,2}\)
split}{C.9 Final E\_\{2,2\} split}}\label{c.9-final-e_22-split}

The two surviving roots are

\[
\begin{aligned}
s^{(0)}={}&(2,2,2,4,6,6,7,7,9,9,9,10,12,12,14,14,15,16,17,17),\\
t^{(0)}={}&229,
\end{aligned}
\tag{C.47}
\]

and

\[
\begin{aligned}
s^{(1)}={}&(2,2,2,5,5,6,6,8,9,9,9,10,12,12,14,14,15,16,17,17),\\
t^{(1)}={}&231.
\end{aligned}
\tag{C.48}
\]

Each has eleven degree classes and hence \(11\cdot19=209\) spectral
variables. The checker verifies by name, not merely by position, that
class-model variable index \(209\) is \(E_{2,2}\). Both sequences have
three degree-\(2\) vertices, so for an actual graph

\[
E_{2,2}\in\{0,1,2,3\}.
\tag{C.49}
\]

For each root, the adjacent children are

\[
E_{2,2}\le0\qquad\text{and}\qquad E_{2,2}\ge1.
\tag{C.50}
\]

The left child fixes \(E_{2,2}=0\); the right contains \(1,2,3\). The
children are disjoint and exhaustive for the actual integer count. Each
of the four child continuous relaxations has an exact Farkas
certificate. Thus

\[
63{,}235+16{,}542+343+4=80{,}124
\tag{C.51}
\]

exact leaves cover all complement representatives, with each final
representative replaced by two branch leaves.

\subsubsection{C.10 Universal graph-to-model
implication}\label{c.10-universal-graph-to-model-implication}

Suppose a target graph existed. Equations (C.1)--(C.5) put its degree
sequence and triangle count in the complete enumeration. After
complementing if necessary, it realizes a retained representative. Its
eigenprojector diagonals give \(Y\) satisfying Stage 1. Its actual
neighbor-degree sums give \(S\) satisfying Stage 2. Its actual edge,
graph-triangle, and complement-triangle counts give \(E,T,C\) satisfying
every Stage-3 equation, capacity, and incidence inequality. If it
reaches a final root, its integer \(E_{2,2}\) lies in exactly one child
of (C.50).

The exact certificate at the corresponding leaf contradicts feasibility
by (C.45), and the verified set partition ensures that no candidate
avoids a leaf. Therefore no target graph exists.

\subsubsection{C.11 Reproduction identities and provenance
note}\label{c.11-reproduction-identities-and-provenance-note}

The release launcher pins the following internal helper files:

\begin{verbatim}
verification/n20/verify_enumeration_v2.py
4df5db86179076364561689d8a4e62cba67890ded650ecc354f53949492bdbd9

verification/n20/verify_certificates_v2.py
b12dfe89fa5b8562844da9f963fdc9782336a3ef29219f6d481dea71170d385e

verification/n20/audit_cycle_identity_v2.py
c59c0255ba40f7b311c614b7493b760460a622ceb09fca08ce1669230248f956
\end{verbatim}

With those helpers, the telemetry-free canonical JSON has SHA-256

\begin{verbatim}
ebc3d8065c3f0c2e869e940f2d50c572292c45646a5ee6ac658a3f805b1ccbac
\end{verbatim}

The combined release retains two fresh full macOS replay transcripts for
this current helper set:

\begin{verbatim}
logs/n20-current-run1.log
9b7ad79994af963e0ba8acfce70d2c3f29c451fb914eae4d57a4d98010989688

logs/n20-current-run2.log
f85d291f862188b2f9b653f512ea458b971f9a4910f79cddf275a2b8f302b3ec
\end{verbatim}

Both runs exit 0 and reproduce the canonical bytes under the offline
sandbox. The expected terminal statuses are:

\begin{verbatim}
INDEPENDENT_ENUMERATION_V2_PASS
INDEPENDENTLY_VERIFIED_UNSAT
INDEPENDENT_REPRODUCTION_PASS
\end{verbatim}

These are internal cross-check status labels; Section 7.3 states their
trust boundary. Earlier wrapper identities are recorded in
\texttt{ARTIFACT\_PROVENANCE.md}.

\subsection{Data and code
availability}\label{data-and-code-availability}

The archival release contains this manuscript and source, the two frozen
certificate archives, their detached hashes, the hardened order-20
wrapper and internal v2 verifiers, the expected canonical order-20 JSON,
exact reproduction programs, and complete replay logs. The top-level
manifest records the SHA-256 digest of every deposited file. The
accompanying code repository is available on
\href{https://github.com/shblue21/s16-s20-laplacian-certified-proofs}{GitHub}.

\subsection{Author contributions and
responsibility}\label{author-contributions-and-responsibility}

Jihun Kim conceived the study, performed and verified the computations,
prepared the deposited materials and wrote the manuscript. He takes
responsibility for the work.

\subsection{AI assistance disclosure}\label{ai-assistance-disclosure}

OpenAI ChatGPT Pro and Codex assisted with research planning, drafting,
code development, test design and internal review. The author checked
the mathematical claims, executed the final verification and approved
the manuscript and deposited files.

\subsection{Funding}\label{funding}

This research received no external funding.

\subsection{Competing interests}\label{competing-interests}

The author declares no competing interests.

\subsection{References}\label{references}

\begingroup\small

\begin{enumerate}
\def\labelenumi{\arabic{enumi}.}
\tightlist
\item
  S. M. Fallat, S. J. Kirkland, J. J. Molitierno and M. Neumann, ``On
  graphs whose Laplacian matrices have distinct integer eigenvalues,''
  \emph{Journal of Graph Theory} \textbf{50} (2005), 162--174.
  \href{https://doi.org/10.1002/jgt.20102}{DOI}.
\item
  K. Ch. Das, S.-G. Lee and G.-S. Cheon, ``On the conjecture for certain
  Laplacian integral spectrum of graphs,'' \emph{Journal of Graph
  Theory} \textbf{63} (2010), 106--113.
  \href{https://doi.org/10.1002/jgt.20412}{DOI}.
\item
  A. Goldberger and M. Neumann, ``On a conjecture on a Laplacian matrix
  with distinct integral spectrum,'' \emph{Journal of Graph Theory}
  \textbf{72} (2013), 178--208.
  \href{https://doi.org/10.1002/jgt.21638}{DOI}.
\item
  A. Hameed, Z. U. Khan and M. Tyaglov, ``Laplacian energy and first
  Zagreb index of Laplacian integral graphs,'' \emph{Analele Științifice
  ale Universității Ovidius Constanța} \textbf{30} (2022), 133--160.
  \href{https://doi.org/10.2478/auom-2022-0023}{DOI}.
\item
  N. Johnston, S. Plosker and L. M. B. Varona, ``Enumeration of
  Laplacian integral and \(\{-1,0,1\}\)-diagonalizable graphs,''
  arXiv:2607.06336v1 (2026).
  \href{https://doi.org/10.48550/arXiv.2607.06336}{DOI}.
\item
  N. Biggs, \emph{Algebraic Graph Theory}, 2nd ed., Cambridge University
  Press,

  \begin{enumerate}
  \def\labelenumii{\arabic{enumii}.}
  \setcounter{enumii}{1992}
  \tightlist
  \item
  \end{enumerate}
\item
  P. Erdős and T. Gallai, ``Graphs with prescribed degrees of
  vertices,'' \emph{Matematikai Lapok} \textbf{11} (1960), 264--274.
\item
  S. L. Hakimi, ``On realizability of a set of integers as degrees of
  the vertices of a linear graph. I,'' \emph{Journal of the Society for
  Industrial and Applied Mathematics} \textbf{10} (1962), 496--506.
\item
  D. J. Kleitman and D. L. Wang, ``Algorithms for constructing graphs
  and digraphs with given valences and factors,'' \emph{Discrete
  Mathematics} \textbf{6} (1973), 79--88.
\item
  A. Schrijver, \emph{Theory of Linear and Integer Programming}, Wiley,
  1986.
\end{enumerate}

\endgroup

\end{document}
